\chapter{角动量}
\label{chap:angular-momentum-tensors}

两个角动量同时存在时，直积基 $|j_1m_1\rangle|j_2m_2\rangle$ 逐个标记两个分量，却不直接显露总转动下的不可约块。总角动量基则对角化 $\vb J^2$ 和 $J_z$，Clebsch--Gordan 系数正是两组基之间的变换矩阵。

算符本身也会在转动下组成不可约张量。把态的几何信息和算符的几何信息分开以后，Wigner--Eckart 定理把全部磁量子数依赖压缩进 Clebsch--Gordan 系数，剩下的约化矩阵元才包含具体动力学。

\section{Clebsch--Gordan 分解与角动量耦合}
\label{sec:cg-coupling}


两个独立角动量放在一起时，总 Hilbert 空间不是两个空间的直接和，而是张量积
\[
\mathcal H=\mathcal H_{j_1}\otimes\mathcal H_{j_2}.
\]
群在两个因子上同时转动，因此总表示是 $D^{(j_1)}\otimes D^{(j_2)}$。这个表示一般可约，所以真正的问题是：它分解成哪些总角动量 $j$ 的不可约表示？

“非耦合基”$|j_1m_1\rangle|j_2m_2\rangle$ 让 $J_{1z},J_{2z}$ 对角；“耦合基”$|JM\rangle$ 让 $J^2,J_z$ 对角。Clebsch--Gordan 系数只是这两组正交基之间的基变换矩阵元：
\[
|JM\rangle=\sum_{m_1,m_2}|j_1m_1;j_2m_2\rangle
\langle j_1m_1,j_2m_2|JM\rangle.
\]
因此 CG 系数并不是一张孤立的数字表，它们来自张量积表示的不可约分解和基变换。


设两个彼此独立的角动量自由度分别承载 $SU(2)$ 不可约表示 $D^{(j_1)}$ 与 $D^{(j_2)}$。总 Hilbert 空间是
\[
\mathcal H_{j_1}\otimes\mathcal H_{j_2},
\qquad
\dim=(2j_1+1)(2j_2+1),
\]
群作用为张量积表示
\[
U(g)=U_{j_1}(g)\otimes U_{j_2}(g).
\]
在这个空间中，自然的非耦合基是
\[
\lvert j_1,m_1\rangle\otimes\lvert j_2,m_2\rangle.
\]
但总生成元
\[
\vb J=\vb J_1+\vb J_2
\]
仍满足同一个 $\mathfrak{su}(2)$ 代数，所以整个张量积表示必须分解成若干 $D^{(j)}$。寻找这种分解，等价于同时对角化 $J^2,J_z$，把非耦合基换成耦合基 $\ket{(j_1j_2)jm}$。

\begin{theorem}[Clebsch--Gordan 分解]
对任意 $j_1,j_2\in\{0,\tfrac12,1,\dots\}$，有
\begin{equation}
D^{(j_1)}\otimes D^{(j_2)}
\cong
\bigoplus_{j=|j_1-j_2|}^{j_1+j_2}D^{(j)},
\label{eq:cg-decomposition}
\end{equation}
其中 $j$ 每次增加 $1$，且每个不可约表示恰出现一次。

\end{theorem}

\begin{proof}
张量积空间中 $J_z=J_{1z}+J_{2z}$，所以最大可能磁量子数为
\[
m_{\max}=j_1+j_2,
\]
且只有唯一态
\[
\lvert j_1,j_1\rangle\lvert j_2,j_2\rangle
\]
具有该 $m$。它必然是某个不可约分量的最高权，因此该分量的 $j=j_1+j_2$。对这个最高权态反复作用总降阶算符
\[
J_-=J_{1-}+J_{2-}
\]
便生成一个完整的 $D^{(j_1+j_2)}$，维数 $2(j_1+j_2)+1$。

从总空间中去掉这个不变子空间后，剩余空间仍是 $SU(2)$ 表示。考察 $m=j_1+j_2-1$ 的子空间：原来有两个非耦合基，其中一个已经属于刚才的最高 $j$ 多重态，正交补只剩一个方向，所以它必是下一个不可约分量的最高权，标签为 $j_1+j_2-1$。继续同样的“按 $m$ 从高到低减去已生成多重态”过程，会依次得到
\[
j_1+j_2,\ j_1+j_2-1,\dots.
\]
何时终止可由维数确定。假设 $j_1\ge j_2$，若一直到 $j=j_1-j_2$，右侧总维数为
\begin{align*}
\sum_{j=j_1-j_2}^{j_1+j_2}(2j+1)
&=(2j_2+1)\frac{[2(j_1-j_2)+1]+[2(j_1+j_2)+1]}2\\
&=(2j_2+1)(2j_1+1),
\end{align*}
恰好等于张量积空间维数。因此没有遗漏，也没有重数大于 $1$ 的分量，得到 \eqnrefs{eq:cg-decomposition}。
\end{proof}

三角规则还可以通过最高权构造理解。张量积中最高磁量子数为 $M_{\max}=j_1+j_2$，而满足这个 $M$ 的态只有
\[
|j_1,j_1\rangle\otimes|j_2,j_2\rangle,
\]
所以它一定是一个总角动量 $J=j_1+j_2$ 不可约表示的最高权态。对它反复作用总降算符 $J_-=J_{1-}+J_{2-}$，得到整个 $J=j_1+j_2$ 多重态。

把这个 $(2J+1)$ 维子空间从张量积空间中正交剔除后，剩余空间若非零，再找其中最高 $M$ 的态，它会成为下一不可约表示 $J=j_1+j_2-1$ 的最高权态。重复这一过程，直到 $J=|j_1-j_2|$。因此
\[
j_1\otimes j_2=\bigoplus_{J=|j_1-j_2|}^{j_1+j_2}J,
\]
且每个 $J$ 恰好出现一次。最后可用维数恒等式
\[
(2j_1+1)(2j_2+1)=\sum_{J=|j_1-j_2|}^{j_1+j_2}(2J+1)
\]
检查没有遗漏任何维度。


这个证明值得注意：三角条件
\[
|j_1-j_2|\le j\le j_1+j_2
\]
不是经典矢量加法的经验类比，而是最高权结构与维数计数共同强制出来的表示论结论。

张量积必须再约化，而约化的核心工具是 Clebsch--Gordan 系数。下面从递推关系出发系统地构造它们。

两套正交基之间的酉变换系数定义为
\begin{equation}
\ket{(j_1j_2)jm}
=
\sum_{m_1,m_2}
C^{jm}_{j_1m_1,j_2m_2}
\lvert j_1,m_1\rangle\lvert j_2,m_2\rangle,
\label{eq:cg-definition}
\end{equation}
其中
\[
C^{jm}_{j_1m_1,j_2m_2}=0
\qquad\text{除非}\qquad m=m_1+m_2.
\]
这些 $C$ 就是 Clebsch--Gordan 系数。因为只是正交基变换，它们满足完备与正交关系
\begin{align}
\sum_{m_1,m_2}
C^{jm}_{j_1m_1,j_2m_2}
C^{j'm'}_{j_1m_1,j_2m_2}{}^*
&=\delta_{jj'}\delta_{mm'},\label{eq:cg-orthogonality-1}\\
\sum_{j,m}
C^{jm}_{j_1m_1,j_2m_2}
C^{jm}_{j_1m_1',j_2m_2'}{}^*
&=\delta_{m_1m_1'}\delta_{m_2m_2'}.
\label{eq:cg-orthogonality-2}
\end{align}

系数不必整张背诵。实际计算可以从最高权态出发，用总升降算符递推。由
\[
J_\pm=J_{1\pm}+J_{2\pm}
\]
作用到 \eqnrefs{eq:cg-definition} 两边并比较同一非耦合基系数，得到
\begin{align}
&\sqrt{(j\mp m)(j\pm m+1)}\,
C^{j,m\pm1}_{j_1m_1,j_2m_2}
\nonumber\\
&\quad=
\sqrt{(j_1\pm m_1)(j_1\mp m_1+1)}
C^{jm}_{j_1,m_1\mp1,j_2m_2}
\nonumber\\
&\qquad+
\sqrt{(j_2\pm m_2)(j_2\mp m_2+1)}
C^{jm}_{j_1m_1,j_2,m_2\mp1}.
\label{eq:cg-recursion}
\end{align}
配合归一化和固定相位约定（通常采用 Condon--Shortley 约定），这就能递归生成全部系数。

\begin{exampleplain}[两个自旋 $1/2$]
由
\[
\frac12\otimes\frac12=1\oplus0
\]
可知四维张量积空间分为三维三重态和一维单态。最高权态唯一：
\[
\ket{1,1}=\ket{\uparrow\uparrow}.
\]
作用 $J_-=J_{1-}+J_{2-}$，
\[
\sqrt2\hbar\ket{1,0}
=\hbar\ket{\downarrow\uparrow}
 +\hbar\ket{\uparrow\downarrow},
\]
所以
\[
\ket{1,0}=\frac{\ket{\uparrow\downarrow}+\ket{\downarrow\uparrow}}{\sqrt2},
\qquad
\ket{1,-1}=\ket{\downarrow\downarrow}.
\]
四维空间中与这三个态正交的唯一归一化组合为
\[
\ket{0,0}
=\frac{\ket{\uparrow\downarrow}-\ket{\downarrow\uparrow}}{\sqrt2}.
\]
这正是\chapref{chap:permutation-group}讨论两电子费米统计时会反复使用的“自旋三重态交换对称、自旋单态交换反对称”。
\end{exampleplain}

\begin{proposition}[$j\otimes\tfrac12$ 的显式耦合]
在 Condon--Shortley 相位约定下，设第一个角动量为 $j$，第二个为 $1/2$，则
\begin{align}
\ket{j+\tfrac12,M}
&=
\sqrt{\frac{j+M+\tfrac12}{2j+1}}
\lvert j,M-\tfrac12\rangle\ket{\tfrac12,\tfrac12}
\nonumber\\
&\quad+
\sqrt{\frac{j-M+\tfrac12}{2j+1}}
\lvert j,M+\tfrac12\rangle\ket{\tfrac12,-\tfrac12},
\label{eq:cg-j-half-plus}\\[2mm]
\ket{j-\tfrac12,M}
&=
-\sqrt{\frac{j-M+\tfrac12}{2j+1}}
\lvert j,M-\tfrac12\rangle\ket{\tfrac12,\tfrac12}
\nonumber\\
&\quad+
\sqrt{\frac{j+M+\tfrac12}{2j+1}}
\lvert j,M+\tfrac12\rangle\ket{\tfrac12,-\tfrac12}.
\label{eq:cg-j-half-minus}
\end{align}
\end{proposition}
\begin{proof}
固定总磁量子数 $M$。除端点处只有一个基矢量外，对应的非耦合子空间由
\[
\ket{+}
=\lvert j,M-\tfrac12\rangle\ket{\tfrac12,\tfrac12},
\qquad
\ket{-}
=\lvert j,M+\tfrac12\rangle\ket{\tfrac12,-\tfrac12}
\]
张成。与其从最高权态逐级递推，不如直接在这个二维空间中对角化 $J^2$。利用
\[
J^2=J_1^2+S^2+2J_{1z}S_z+J_{1+}S_-+J_{1-}S_+,
\]
其中第二个角动量记为 $\vb S$，且 $s=1/2$。把恒定部分 $j(j+1)\hbar^2+\frac34\hbar^2$ 提出，定义
\[
A=\frac{J^2-j(j+1)\hbar^2-\frac34\hbar^2}{\hbar^2}.
\]
先算对角项。由于 $S_z\ket{\uparrow}=\frac\hbar2\ket{\uparrow}$、$S_z\ket{\downarrow}=-\frac\hbar2\ket{\downarrow}$，有
\[
\frac{2J_{1z}S_z}{\hbar^2}\ket{+}
=\left(M-\frac12\right)\ket{+},
\qquad
\frac{2J_{1z}S_z}{\hbar^2}\ket{-}
=-\left(M+\frac12\right)\ket{-}.
\]
再算非对角项。对自旋 $1/2$，
\[
S_-\ket{\uparrow}=\hbar\ket{\downarrow},
\qquad
S_+\ket{\downarrow}=\hbar\ket{\uparrow},
\]
而
\[
J_{1+}\lvert j,M-\tfrac12\rangle
=\hbar\sqrt{\left(j+\frac12\right)^2-M^2}\,
\lvert j,M+\tfrac12\rangle,
\]
\[
J_{1-}\lvert j,M+\tfrac12\rangle
=\hbar\sqrt{\left(j+\frac12\right)^2-M^2}\,
\lvert j,M-\tfrac12\rangle.
\]
因此在有序基 $(\ket{+},\ket{-})$ 中
\begin{equation}
[A]_M=
\begin{pmatrix}
M-\frac12 & C_M\\
C_M & -M-\frac12
\end{pmatrix},
\qquad
C_M=\sqrt{\left(j+\frac12\right)^2-M^2}.
\label{eq:j-half-two-by-two}
\end{equation}
由 Clebsch--Gordan 分解，这个二维空间中的两个本征态分别属于 $J=j+\frac12$ 与 $J=j-\frac12$。对前者，
\[
J(J+1)-j(j+1)-\frac34=j,
\]
所以 $A$ 的对应本征值为 $j$。设本征向量为 $a\ket{+}+b\ket{-}$，则
\[
\left(M-j-\frac12\right)a+C_M b=0.
\]
利用
\[
C_M^2
=\left(j+M+\frac12\right)\left(j-M+\frac12\right)
\]
可写成
\[
\frac{b}{a}
=\sqrt{\frac{j-M+\frac12}{j+M+\frac12}}.
\]
归一化 $|a|^2+|b|^2=1$，并按 Condon--Shortley 约定取最高权系数为正，得到
\[
a=\sqrt{\frac{j+M+\frac12}{2j+1}},
\qquad
b=\sqrt{\frac{j-M+\frac12}{2j+1}},
\]
即 \eqnrefs{eq:cg-j-half-plus}。
剩余的归一化方向必须与它正交。取
\[
\binom{a_-}{b_-}
=
\binom{-b}{a},
\]
便有
\[
[A]_M\binom{-b}{a}=-(j+1)\binom{-b}{a}.
\]
而
\[
\left(j-\frac12\right)\left(j+\frac12\right)
-j(j+1)-\frac34=-(j+1),
\]
所以该向量正属于 $J=j-\frac12$ 多重态。代入 $a,b$ 就得到 \eqnrefs{eq:cg-j-half-minus}。端点处某个平方根自动为零，因此同一公式也覆盖只有一个非耦合基矢量的情形。
\end{proof}
前面建立了 $SO(3)$ 的全部不可约表示 $D^{(j)}$。现在把它直接用到分子转动光谱上——这是表示论给出物理能级的最干净例子之一。
\medskip
\noindent\emph{刚性转子的能级。}
刚性转子（如双原子分子或球形陀螺）的哈密顿量为
\[
\hat H=\frac{\hat{\vb J}^2}{2I},
\]
其中 $I$ 为转动惯量。$\hat{\vb J}^2$ 的本征值已由角动量代数给出：
\[
\hat{\vb J}^2\ket{j,m}=\hbar^2 j(j+1)\ket{j,m}.
\]
因此能级为
\[
E_j=B\,j(j+1),
\qquad
B=\frac{\hbar^2}{2I},
\]
每个能级 $E_j$ 的简并度为 $2j+1$（$m=-j,\dots,j$）。
$\hat H$ 与所有转动对易，而 $\hat{\vb J}^2$ 在不可约表示 $D^{(j)}$ 上由
Schur 引理固定为 $\hbar^2j(j+1)$。因此每个 $j$ 多重态正好包含
$\dim D^{(j)}=2j+1$ 个磁量子数态；这一简并度无需求解转子微分方程。
\par\medskip
Clebsch--Gordan 系数 $C(j_1,j_2,j_3;m_1,m_2,m_3)$ 在物理中极为常用，但它对三个角动量的处理不对称（$j_3$ 被特殊对待）。Wigner 引入的 $3j$ 符号通过一个相位因子把 CG 系数对称化，使得三个角动量地位平等。
\begin{definitionplain}[Wigner 3j 符号]
\[
\begin{pmatrix}j_1&j_2&j_3\\m_1&m_2&m_3\end{pmatrix}
=\frac{(-1)^{j_1-j_2-m_3}}{\sqrt{2j_3+1}}\,
C(j_1,j_2,j_3;m_1,m_2,-m_3).
\]
\label{def:3j-symbol}
\end{definitionplain}
\begin{proposition}[3j 符号的对称性质]
$3j$ 符号满足：
\begin{enumerate}[label=(\roman*),leftmargin=3em]
\item \textbf{列置换偶交换不变：}
\[
\begin{pmatrix}j_2&j_3&j_1\\m_2&m_3&m_1\end{pmatrix}
=
\begin{pmatrix}j_1&j_2&j_3\\m_1&m_2&m_3\end{pmatrix}.
\]
\item \textbf{列置换奇交换出相位：}
\[
\begin{pmatrix}j_2&j_1&j_3\\m_2&m_1&m_3\end{pmatrix}
=(-1)^{j_1+j_2+j_3}
\begin{pmatrix}j_1&j_2&j_3\\m_1&m_2&m_3\end{pmatrix}.
\]
\item \textbf{磁量子数反号：}
\[
\begin{pmatrix}j_1&j_2&j_3\\-m_1&-m_2&-m_3\end{pmatrix}
=(-1)^{j_1+j_2+j_3}
\begin{pmatrix}j_1&j_2&j_3\\m_1&m_2&m_3\end{pmatrix}.
\]
\item \textbf{选择定则：}非零要求 $m_1+m_2+m_3=0$ 且 $|j_i-j_k|\le j_l\le j_i+j_k$（三角形条件）。
\end{enumerate}
\label{prop:3j-symmetry}
\end{proposition}
$3j$ 符号直接给出 Wigner--Eckart 分解\cref{thm:wigner-eckart-angular} 中的几何因子
\[
(-1)^{j'-m'}\begin{pmatrix}j'&k&j\\-m'&q&m\end{pmatrix}.
\]
同一个对象也控制三个球谐函数的 Gaunt 积分：
\[
\int Y_{\ell_1 m_1}^*Y_{\ell_2 m_2}Y_{\ell_3 m_3}\,d\Omega
=\sqrt{\frac{(2\ell_1+1)(2\ell_2+1)(2\ell_3+1)}{4\pi}}
\begin{pmatrix}\ell_1&\ell_2&\ell_3\\0&0&0\end{pmatrix}
\begin{pmatrix}\ell_1&\ell_2&\ell_3\\-m_1&m_2&m_3\end{pmatrix}.
\]
因此 $3j$ 符号没有引入新的耦合数据；它把 CG 系数乘上相位与归一化因子，换来三个角动量标签的对称地位，这在再耦合计算中尤其方便。
前面通过递推关系确定了 CG 系数。现在给出 Racah 1942 年的闭合表达式，它把 CG 系数写成阶乘的显式组合，无需递推即可直接计算任意 $C(j_1,j_2,j_3;m_1,m_2,m_3)$。
\begin{theorem}[Racah 闭合公式]
\[
\begin{aligned}
C(j_1,j_2,j_3;m_1,m_2,m_3)
&=\delta_{m_1+m_2,m_3}\,\Delta(j_1,j_2,j_3)
\sum_z(-1)^z\\
&\quad\times
\frac{
\sqrt{\begin{gathered}
(j_1+m_1)!(j_1-m_1)!(j_2+m_2)!\\[-2pt]
{}\times(j_2-m_2)!(j_3+m_3)!(j_3-m_3)!
\end{gathered}}
}{
\begin{gathered}
(j_1+j_2-j_3-z)!(j_1-m_1-z)!(j_2+m_2-z)!\\[-2pt]
{}\times(j_3-j_2+m_1+z)!(j_3-j_1-m_2+z)!\,z!
\end{gathered}
}.
\end{aligned}
\]
其中三角形因子
\[
\Delta(j_1,j_2,j_3)=\sqrt{\frac{(j_1+j_2-j_3)!(j_1-j_2+j_3)!(-j_1+j_2+j_3)!}{(j_1+j_2+j_3+1)!}},
\]
求和遍历所有使阶乘自变量为非负整数的 $z$。
\label{thm:racah-cg}
\end{theorem}
\par\medskip
\noindent\textbf{Racah 公式的推导}\quad
先把最高权态归一化：从 $m_1=j_1,m_2=j_2,m_3=j_1+j_2$ 的唯一态出发：
\[
\ket{j_1 j_1;j_2 j_2;j_3,j_1+j_2}=\ket{j_1 j_1}\ket{j_2 j_2},
\qquad
C(j_1,j_2,j_1+j_2;j_1,j_2,j_1+j_2)=1.
\]
这固定了整体相位约定（Condon--Shorty 约定）。
随后用降阶算符递推：对最高权态作用 $J_-=J_{1-}+J_{2-}$ 共 $n=j_1+j_2-m_3$ 次。左边给出
\[
J_-^n\ket{j_3,j_3}
=\sqrt{\frac{(2j_3)!\,n!}{(2j_3-n)!}}\ket{j_3,j_3-n}
=\sqrt{\frac{(2j_3)!\,(j_1+j_2-m_3)!}{(j_1+j_2+m_3)!}}\ket{j_3,m_3}.
\]
右边用二项式定理展开 $(J_{1-}+J_{2-})^n$：
\[
(J_{1-}+J_{2-})^n=\sum_{n_1=0}^{n}\binom{n}{n_1}J_{1-}^{n_1}J_{2-}^{n-n_1},
\]
每项含
\[
J_{1-}^{n_1}\ket{j_1 j_1}
=\sqrt{\frac{(2j_1)!\,n_1!}{(2j_1-n_1)!}}\ket{j_1,j_1-n_1},
\quad
J_{2-}^{n-n_1}\ket{j_2 j_2}
=\sqrt{\frac{(2j_2)!\,(n-n_1)!}{(2j_2-n+n_1)!}}\ket{j_2,j_2-n+n_1}.
\]
令 $m_1=j_1-n_1$，$m_2=j_2-(n-n_1)=m_3-m_1$，提取 CG 系数的递推关系：
\[
\begin{aligned}
C(j_1,j_2,j_3;m_1,m_2,m_3)
&=\sqrt{\frac{(2j_3)!\,(j_1+j_2-m_3)!}{(j_1+j_2+m_3)!}}\\
&\quad\times\sum_{n_1}\binom{n}{n_1}
\sqrt{\frac{(2j_1)!\,n_1!\,(2j_2)!\,(n-n_1)!}
{(2j_1-n_1)!\,(2j_2-n+n_1)!}}\,(\text{相位}).
\end{aligned}
\]
接着解递推关系得到闭合式：递推关系是超几何型的。设
\[
\begin{aligned}
C(j_1,j_2,j_3;m_1,m_2,m_3)
&=N\sum_z(-1)^z
\frac{1}{
\begin{gathered}
z!\,(j_1+j_2-j_3-z)!\,(j_1-m_1-z)!\\[-2pt]
{}\times(j_2+m_2-z)!(j_3-j_2+m_1+z)!
(j_3-j_1-m_2+z)!
\end{gathered}}.
\end{aligned}
\]
代入递推关系，利用 Vandermonde 恒等式 $\sum_k\binom{a}{k}\binom{b}{n-k}=\binom{a+b}{n}$ 化简，确认求和满足 $_3F_2$ 级数的终止条件（某个分子阶乘自变量取负整数时级数截断），求和范围由所有阶乘自变量非负确定。归一化条件 $\sum_{m_1,m_2}|C|^2=1$ 确定常数 $N$：利用 $\sum_{m_3}|C|^2=1$ 对每个 $m_3$ 求和，交换求和顺序后用 Vandermonde 恒等式，得
\[
\begin{aligned}
N={}&\Delta(j_1,j_2,j_3)
\Bigl[(j_1+j_2-m_3)!\,(j_1-j_2+m_3)!\,
(-j_1+j_2+m_3)!\\
&\quad\times(j_1+m_1)!\,(j_1-m_1)!\,(j_2+m_2)!\,(j_2-m_2)!\\
&\quad\times(j_3+m_3)!\,(j_3-m_3)!\Bigr]^{1/2}.
\end{aligned}
\]
交换对称性 $C(j_2,j_1,j_3;-m_2,-m_1,-m_3)=(-1)^{j_1+j_2-j_3}C(j_1,j_2,j_3;m_1,m_2,m_3)$ 确定相位约定，与 Condon--Shorty 约定一致。
\par\medskip
\medskip
\noindent\emph{用 Racah 公式计算 $C(1,\frac12,\frac32;0,\frac12,\frac12)$。}
取 $j_1=1,j_2=\frac12,j_3=\frac32,m_1=0,m_2=\frac12,m_3=\frac12$。三角形因子：
\[
\Delta(1,\tfrac12,\tfrac32)=\sqrt{\frac{0!\,2!\,1!}{4!}}=\sqrt{\frac{2}{24}}=\frac1{\sqrt{12}}.
\]
分子平方根：
\[
\sqrt{1!\,1!\,1!\,0!\,2!\,1!}=\sqrt{2}.
\]
求和变量 $z$ 使所有阶乘自变量非负：$z\ge0$，$z\le j_1+j_2-j_3=0$，$z\le j_1-m_1=1$，$z\le j_2+m_2=1$，$z\ge j_2-j_3-m_1=-1$，$z\ge j_1-j_3+m_2=0$。唯一 $z=0$。故
\[
C=\frac1{\sqrt{12}}\cdot\sqrt{2}\cdot\frac{1}{0!\,1!\,1!\,0!\,0!\,0!}=\frac{\sqrt{2}}{\sqrt{12}}=\frac1{\sqrt{6}}.
\]
验证：$j=3/2$ 四重态中 $m=1/2$ 态为 $\frac{1}{\sqrt3}\alpha\alpha\beta+\sqrt{\frac23}\beta\alpha\alpha$，$C(1,\frac12,\frac32;0,\frac12,\frac12)$ 对应 $\alpha\beta$ 分量，系数 $\frac1{\sqrt3}\cdot\frac1{\sqrt2}=\frac1{\sqrt6}$（含自旋归一化），吻合。
\par\medskip
三个角动量可按两种方式耦合：
\[
\bigl[(j_1\otimes j_2)\otimes j_3\bigr]_{J}
\qquad\text{vs}\qquad
\bigl[j_1\otimes(j_2\otimes j_3)\bigr]_{J}.
\]
两种基之间的变换矩阵元素就是再耦合系数。6j 符号是其对称化版本。
\begin{definitionplain}[Racah 系数与 6j 符号]
再耦合系数定义为
\[
\braket{(j_1 j_2)J_{12},j_3;JM}{j_1,(j_2 j_3)J_{23};JM}
=\sqrt{(2J_{12}+1)(2J_{23}+1)}\,(-1)^{j_1+j_2+j_3+J}
\begin{Bmatrix}j_1&j_2&J_{12}\\j_3&J&J_{23}\end{Bmatrix},
\]
其中 $\begin{Bmatrix}j_1&j_2&j_3\\j_4&j_5&j_6\end{Bmatrix}$ 是 Wigner 6j 符号。
\end{definitionplain}
\begin{theorem}[6j 符号的 Racah 闭合表达式]
\[
\begin{Bmatrix}j_1&j_2&j_3\\j_4&j_5&j_6\end{Bmatrix}
=\Delta(j_1,j_2,j_3)\Delta(j_1,j_5,j_6)\Delta(j_4,j_2,j_6)\Delta(j_4,j_5,j_3)
\sum_z\frac{(-1)^z(z+1)!}{\prod_{i=1}^4(z-a_i)!\,\prod_{i=1}^3(b_i-z)!},
\]
其中四个 $\Delta$ 是三角形因子，$a_i$ 为三组 $j$ 之和：
\[
a_1=j_1+j_2+j_3,\ a_2=j_1+j_5+j_6,\ a_3=j_4+j_2+j_6,\ a_4=j_4+j_5+j_3,
\]
$b_i$ 为三组 $j$ 之和：
\[
b_1=j_1+j_2+j_4+j_5,\ b_2=j_2+j_3+j_5+j_6,\ b_3=j_3+j_1+j_6+j_4,
\]
求和遍历所有使阶乘自变量为非负整数的 $z$。
\label{thm:6j-racah}
\end{theorem}
\begin{proposition}[6j 符号的对称性质]
6j 符号满足：
\begin{enumerate}[label=(\roman*),leftmargin=3em]
\item \textbf{列置换不变：}交换任意两列不改变值。
\item \textbf{上下行交换不变：}同时对调任意两列中的上下元素不改变值。
\item 以上两类操作生成 24 种等价排列，6j 符号在这 24 种置换下完全不变。
\item \textbf{三角形条件：}$(j_1,j_2,j_3)$、$(j_1,j_5,j_6)$、$(j_4,j_2,j_6)$、$(j_4,j_5,j_3)$ 都必须满足三角形条件。
\end{enumerate}
\end{proposition}
\par\medskip
\noindent\textbf{6j 符号对称性的来源}\quad
把 6j 符号写成四个 3j 符号的磁量子数缩并。交换任意两列时，每个 3j 符号都按相同的 Condon--Shortley 规则取得相位，而每条内部边出现两次，故总相位平方为 $1$。同时交换任意两列的上下元素等价于重新命名内部磁量子数，缩并值不变。
这两类操作把六个角动量视为四面体的六条棱：允许的置换恰为四面体四个顶点的全排列 $S_4$，共有 $24$ 个。四个顶点处的三角形条件正对应定理中列出的四组三角形。因而 6j 的对称性不是额外假设，而是缩并图在重新标记顶点下的不变性。
\par\medskip
几何解释还能控制大角动量极限。若以 $\ell_i=j_i+\tfrac12$ 为六条棱的 Euclidean 四面体存在，Ponzano--Regge 渐近式为
\[
\begin{Bmatrix}j_1&j_2&j_3\\j_4&j_5&j_6\end{Bmatrix}
\sim\frac{1}{\sqrt{12\pi V}}
\cos\left(\sum_{i=1}^{6}\ell_i\theta_i+\frac{\pi}{4}\right),
\]
其中 $V$ 是四面体体积，$\theta_i$ 是与棱 $\ell_i$ 对应的外二面角。相位是 Regge 作用量而不是体积本身。
\medskip
\noindent\emph{$\begin{Bmatrix}j_1&j_2&j_3\\0&j_3&j_2\end{Bmatrix}$ 的简化。}
当一个 $j$ 为零时，6j 符号大幅简化。取 $j_4=0$，三角形条件要求 $j_5=j_3$，$j_6=j_2$。代入 Racah 公式，求和只有一项，得
\[
\begin{Bmatrix}j_1&j_2&j_3\\0&j_3&j_2\end{Bmatrix}
=\frac{(-1)^{j_1+j_2+j_3}}{\sqrt{(2j_2+1)(2j_3+1)}}.
\label{eq:6j-zero}
\]
物理含义：当一个子系统角动量为零，再耦合是平凡的，6j 符号退化为 CG 系数的归一化因子。
\par\medskip
\begin{proposition}[6j 符号正交关系]
\[
\sum_{j_3}(2j_3+1)(2j_6+1)
\begin{Bmatrix}j_1&j_2&j_3\\j_4&j_5&j_6\end{Bmatrix}
\begin{Bmatrix}j_1&j_2&j_3\\j_4&j_5&j_6'\end{Bmatrix}
=\delta_{j_6 j_6'}.
\]
\end{proposition}
\par\medskip
\noindent\textbf{正交关系的证明}\quad
先回忆再耦合矩阵的定义：固定 $j_1,j_2,j_3,J$，再耦合矩阵
\[
U_{J_{12},J_{23}}=\braket{(j_1 j_2)J_{12},j_3;J}{j_1,(j_2 j_3)J_{23};J}
\]
是同一 Hilbert 空间 $\mathcal H_{j_1}\otimes\mathcal H_{j_2}\otimes\mathcal H_{j_3}$ 中两组正交基之间的变换矩阵。第一组基先耦合 $j_1\otimes j_2\to J_{12}$ 再耦合 $J_{12}\otimes j_3\to J$；第二组基先耦合 $j_2\otimes j_3\to J_{23}$ 再耦合 $j_1\otimes J_{23}\to J$。
随后正交性。因两组基都是正交归一的，$U$ 是正交矩阵：
\[
\sum_{J_{12}}U_{J_{12},J_{23}}U_{J_{12},J_{23}'}=\delta_{J_{23},J_{23}'},
\]
\[
\sum_{J_{23}}U_{J_{12},J_{23}}U_{J_{12}',J_{23}}=\delta_{J_{12},J_{12}'}.
\]
接着6j 符号与再耦合矩阵的关系。再耦合矩阵与 6j 符号的关系为
\[
U_{J_{12},J_{23}}=(-1)^{j_1+j_2+j_3+J}\sqrt{(2J_{12}+1)(2J_{23}+1)}
\begin{Bmatrix}j_1&j_2&J_{12}\\j_3&J&J_{23}\end{Bmatrix}.
\]
再代入并化简。将上式代入正交关系：
\[
\sum_{J_{12}}(2J_{12}+1)(2J_{23}+1)
\begin{Bmatrix}j_1&j_2&J_{12}\\j_3&J&J_{23}\end{Bmatrix}
\begin{Bmatrix}j_1&j_2&J_{12}\\j_3&J&J_{23}'\end{Bmatrix}
=\delta_{J_{23},J_{23}'}.
\]
两边除以 $(2J_{23}+1)$（设 $J_{23}'=J_{23}$ 时）或直接整理：
\[
\sum_{J_{12}}(2J_{12}+1)
\begin{Bmatrix}j_1&j_2&J_{12}\\j_3&J&J_{23}\end{Bmatrix}
\begin{Bmatrix}j_1&j_2&J_{12}\\j_3&J&J_{23}'\end{Bmatrix}
=\frac{\delta_{J_{23},J_{23}'}}{2J_{23}+1}.
\]
这就是 6j 符号的正交关系。相位因子 $(-1)^{2(j_1+j_2+j_3+J)}=1$（因 $2\times$ 半整数=整数）在化简中消失。
这说明，正交关系保证不同再耦合方案之间的变换是幺正的——物理结果不依赖于耦合顺序。这在原子物理的 $LS$--$jj$ 耦合变换中直接使用：变换矩阵的幺正性保证能级计算与耦合方案无关。
\par\medskip
对两个电子，$LS$ 与 $jj$ 耦合基之间的幺正变换由 9j 符号给出：
\[
\begin{aligned}
&\langle(l_1s_1)j_1,(l_2s_2)j_2;J
\mid(l_1l_2)L,(s_1s_2)S;J\rangle\\
&\quad=\sqrt{(2j_1+1)(2j_2+1)(2L+1)(2S+1)}
\begin{Bmatrix}l_1&l_2&L\\s_1&s_2&S\\j_1&j_2&J\end{Bmatrix}.
\end{aligned}
\]
6j 符号是三个角动量改变耦合顺序的矩阵元；Racah $W$ 系数只是它乘以约定相位后的记法。它因而自然出现在复合张量算符和二体相互作用的约化矩阵元中。两个电子的 $LS$ 与 $jj$ 基同时重组四个角动量，所需的完整变换则由下一式的 9j 符号给出。
\begin{definitionplain}[9j 符号]
四个角动量有两种再耦合方案：
\[
\bigl[((j_1\otimes j_2)J_{12}\otimes(j_3\otimes j_4)J_{34})J\bigr]
\quad\to\quad
\bigl[((j_1\otimes j_3)J_{13}\otimes(j_2\otimes j_4)J_{24})J\bigr].
\]
变换矩阵元素定义为
\[
\begin{aligned}
&\braket{
\substack{(j_1j_2)J_{12}\\(j_3j_4)J_{34}};J
}{
\substack{(j_1j_3)J_{13}\\(j_2j_4)J_{24}};J
}\\
&\quad=
\sqrt{(2J_{12}+1)(2J_{34}+1)(2J_{13}+1)(2J_{24}+1)}
\begin{Bmatrix}j_1&j_2&J_{12}\\j_3&j_4&J_{34}\\J_{13}&J_{24}&J\end{Bmatrix}.
\end{aligned}
\]
\end{definitionplain}
\begin{proposition}[9j 符号与 6j 符号的关系]
9j 符号可展开为 6j 符号之和：
\[
\begin{Bmatrix}j_1&j_2&J_{12}\\j_3&j_4&J_{34}\\J_{13}&J_{24}&J\end{Bmatrix}
=\sum_k(-1)^{2k}(2k+1)
\begin{Bmatrix}j_1&j_3&J_{13}\\j_4&J&k\end{Bmatrix}
\begin{Bmatrix}j_2&j_4&J_{24}\\j_3&k&J_{34}\end{Bmatrix}
\begin{Bmatrix}J_{12}&J_{34}&J\\k&j_1&j_2\end{Bmatrix}.
\]
\end{proposition}
\begin{proposition}[9j 符号的对称性质]
9j 符号满足：
\begin{enumerate}[label=(\roman*),leftmargin=3em]
\item \textbf{行置换：}偶置换不变，奇置换出 $(-1)^{\Sigma}$，其中 $\Sigma$ 为所有 9 个 $j$ 之和。
\item \textbf{列置换：}同上。
\item \textbf{行列转置：}矩阵转置不变。
\item 以上共 72 种对称操作。
\end{enumerate}
\end{proposition}
\medskip
\noindent\emph{含零的 9j 符号简化。}
当 $J_{34}=0$（即 $j_3=j_4$，$J_{12}=J$），9j 符号退化为 6j：
\[
\begin{Bmatrix}j_1&j_2&J\\j_3&j_3&0\\J_{13}&J_{24}&J\end{Bmatrix}
=\frac{(-1)^{j_2+j_3+J+J_{13}}}{\sqrt{(2J+1)(2j_3+1)}}
\begin{Bmatrix}j_1&j_2&J\\J_{24}&J_{13}&j_3\end{Bmatrix}.
\]
\par\medskip
\medskip
\noindent 9j 符号的一个重要应用，是联系两个电子的 \(LS\) 耦合与 \(jj\) 耦合。相应基矢之间的变换为
\[
\braket{(l_1 l_2)L,(s_1 s_2)S;J}{(l_1 s_1)j_1,(l_2 s_2)j_2;J}
=\sqrt{(2L+1)(2S+1)(2j_1+1)(2j_2+1)}
\begin{Bmatrix}l_1&l_2&L\\s_1&s_2&S\\j_1&j_2&J\end{Bmatrix}.
\]
对电子 $s_1=s_2=\frac12$，这个 9j 符号直接给出从 $LS$ 项符号到 $jj$ 耦合的变换系数，在自旋--轨道耦合不可忽略时（重原子）用于将 $LS$ 好量子数态变换为 $jj$ 好量子数态。
\par\medskip
这套记号按问题层级逐步增加：CG 系数耦合两个角动量，3j 符号把同一数据写成对称形式，6j 符号改变三个角动量的耦合顺序，9j 符号则比较四个角动量的两种成对耦合。Racah 公式的有限求和范围由分母阶乘非负直接决定，归一化与相位分别由 CG 正交关系和 Condon--Shortley 约定固定。
把同一递推改写成终止的 \({}_3F_2(1)\) 只是一种等价求和方法，不再另设重复推导框。
对常用小角动量值，Racah 公式给出完全显式的 CG 系数表。下面以 $j_1=1\otimes j_2=\frac12$ 为例，展示从最高权态出发、通过降阶算符逐步构造全部 CG 系数的完整过程。物理上，这对应轨道角动量 $l=1$（$p$ 电子）与自旋 $s=\frac12$ 的耦合，给出 $j=\frac12$（$p_{1/2}$）和 $j=\frac32$（$p_{3/2}$）两个多重态——正是自旋--轨道耦合劈裂的群论来源。
\begin{table}[htbp]
\centering
\caption{$j_1=1,\ j_2=\frac12$ 的 CG 系数表（$j_3=\frac12,\frac32$）}
\label{tab:cg-1-half}
\small
\begin{tabular}{c|cc|cc}
\toprule
& \multicolumn{2}{c|}{$j_3=\frac32$} & \multicolumn{2}{c}{$j_3=\frac12$} \\
$m_1,m_2$ & $m_3=\frac32$ & $m_3=\frac12$ & $m_3=\frac12$ & $m_3=-\frac12$ \\
\midrule
$1,\frac12$ & $1$ & $0$ & $0$ & $0$ \\
$1,-\frac12$ & $0$ & $\frac1{\sqrt3}$ & $\sqrt{\frac23}$ & $0$ \\
$0,\frac12$ & $0$ & $\sqrt{\frac23}$ & $-\frac1{\sqrt3}$ & $0$ \\
$0,-\frac12$ & $0$ & $0$ & $0$ & $\sqrt{\frac23}$ \\
$-1,\frac12$ & $0$ & $0$ & $0$ & $\frac1{\sqrt3}$ \\
$-1,-\frac12$ & $0$ & $0$ & $0$ & $1$ \\
\bottomrule
\end{tabular}
\end{table}
\par\medskip
\noindent\textbf{$j_1=1,j_2=\frac12$ CG 表中关键系数的推导}\quad
$C(1,\frac12,\frac32;1,\frac12,\frac32)=1$。最高权态，由 Condon--Shorty 约定直接给出。
$C(1,\frac12,\frac32;0,\frac12,\frac12)=\sqrt{2/3}$。对 $\ket{\frac32,\frac32}=\ket{1,1}\ket{\frac12,\frac12}$ 作用 $J_-$：
\[
J_-\ket{\frac32,\frac32}=\sqrt{j_3(j_3+1)-m_3(m_3-1)}\ket{\frac32,\frac12}=\sqrt{3}\ket{\frac32,\frac12}.
\]
右边 $J_-=J_{1-}+J_{2-}$ 分别作用：
\[
J_{1-}\ket{1,1}\ket{\tfrac12,\tfrac12}=\sqrt{1(2)-1\cdot0}\,\ket{1,0}\ket{\tfrac12,\tfrac12}=\sqrt{2}\ket{1,0}\ket{\tfrac12,\tfrac12},
\]
\[
J_{2-}\ket{1,1}\ket{\tfrac12,\tfrac12}=\sqrt{\tfrac12(\tfrac32)-\tfrac12(-\tfrac12)}\,\ket{1,1}\ket{\tfrac12,-\tfrac12}=\ket{1,1}\ket{\tfrac12,-\tfrac12}.
\]
故
\[
\sqrt{3}\ket{\tfrac32,\tfrac12}=\sqrt{2}\ket{1,0}\ket{\tfrac12,\tfrac12}+\ket{1,1}\ket{\tfrac12,-\tfrac12},
\]
\[
\ket{\tfrac32,\tfrac12}=\sqrt{\tfrac23}\ket{1,0}\ket{\tfrac12,\tfrac12}+\tfrac1{\sqrt3}\ket{1,1}\ket{\tfrac12,-\tfrac12}.
\]
读出 $C(1,\frac12,\frac32;0,\frac12,\frac12)=\sqrt{2/3}$，$C(1,\frac12,\frac32;1,-\frac12,\frac12)=1/\sqrt3$。验证归一化：$\frac{2}{3}+\frac{1}{3}=1$。
$C(1,\frac12,\frac12;1,-\frac12,\frac12)=\sqrt{2/3}$。$j_3=\frac12$ 的态与 $j_3=\frac32$ 正交。设
\[
\ket{\tfrac12,\tfrac12}=a\ket{1,1}\ket{\tfrac12,-\tfrac12}+b\ket{1,0}\ket{\tfrac12,\tfrac12}.
\]
正交条件 $\braket{\frac32,\frac12}{\frac12,\frac12}=0$ 给出
\[
\sqrt{\tfrac23}\,b+\tfrac1{\sqrt3}\,a=0\quad\Longrightarrow\quad a=-\sqrt2\,b.
\]
归一化 $a^2+b^2=1$ 给出 $2b^2+b^2=1$，即 $b=\frac1{\sqrt3}$，$a=-\sqrt{\frac23}$。由 Condon--Shorty 相位约定 $C(1,\frac12,\frac12;1,-\frac12,\frac12)>0$，取 $a=\sqrt{\frac23}$，$b=-\frac1{\sqrt3}$（即整体乘 $-1$）。
$C(1,\frac12,\frac12;0,\frac12,-\frac12)=-1/\sqrt3$。对 $\ket{\frac12,\frac12}=\sqrt{\frac23}\ket{1,1}\ket{\frac12,-\tfrac12}-\frac1{\sqrt3}\ket{1,0}\ket{\tfrac12,\tfrac12}$ 作用 $J_-$：
\[
J_-\ket{\tfrac12,\tfrac12}=\ket{\tfrac12,-\tfrac12}.
\]
右边：
\[
\sqrt{\tfrac23}\Big(\ket{1,0}\ket{\tfrac12,-\tfrac12}+\ket{1,1}\ket{\tfrac12,-\tfrac32}\Big)
-\tfrac1{\sqrt3}\Big(\sqrt2\,\ket{1,-1}\ket{\tfrac12,\tfrac12}+\ket{1,0}\ket{\tfrac12,-\tfrac12}\Big).
\]
但 $j_2=\frac12$ 故 $m_2=-\frac32$ 不存在，该项为零。整理得
\[
\ket{\tfrac12,-\tfrac12}=\Big(\sqrt{\tfrac23}-\tfrac1{\sqrt3}\Big)\ket{1,0}\ket{\tfrac12,-\tfrac12}-\sqrt{\tfrac23}\,\ket{1,-1}\ket{\tfrac12,\tfrac12},
\]
验证系数自洽（$\sqrt{2/3}-1/\sqrt3=1/\sqrt3$），读出 $C(1,\frac12,\frac12;0,\frac12,-\frac12)=-1/\sqrt3$（注意 $m_3=-\frac12$ 分量系数），与表中一致。
\par\medskip
\begin{table}[htbp]
\centering
\caption{$j_1=\frac12,\ j_2=\frac12$ 的 CG 系数表（$j_3=0,1$）}
\label{tab:cg-half-half}
\begin{tabular}{c|cc|cc}
\toprule
& \multicolumn{2}{c|}{$j_3=1$} & \multicolumn{2}{c}{$j_3=0$} \\
$m_1,m_2$ & $m_3=1$ & $m_3=0$ & $m_3=0$ & $m_3=-1$ \\
\midrule
$\frac12,\frac12$ & $1$ & $0$ & $0$ & $0$ \\
$\frac12,-\frac12$ & $0$ & $\frac1{\sqrt2}$ & $\frac1{\sqrt2}$ & $0$ \\
$-\frac12,\frac12$ & $0$ & $\frac1{\sqrt2}$ & $-\frac1{\sqrt2}$ & $0$ \\
$-\frac12,-\frac12$ & $0$ & $0$ & $0$ & $1$ \\
\bottomrule
\end{tabular}
\end{table}
6j 符号满足一个重要的求和恒等式——Biedenharn--Elliott 恒等式，它是再耦合矩阵“结合律”的体现。
\begin{theorem}[Biedenharn--Elliott 恒等式]
\[
\sum_x(-1)^{R+x}(2x+1)
\begin{Bmatrix}j_1&j_2&x\\j_3&j_4&j_5\end{Bmatrix}
\begin{Bmatrix}j_3&j_4&x\\j_6&j_7&j_8\end{Bmatrix}
\begin{Bmatrix}j_6&j_7&x\\j_2&j_1&j_9\end{Bmatrix}
=
\begin{Bmatrix}j_5&j_8&j_9\\j_7&j_1&j_4\end{Bmatrix}
\begin{Bmatrix}j_5&j_8&j_9\\j_6&j_2&j_3\end{Bmatrix},
\]
其中 $R=\sum_{i=1}^9 j_i$，求和遍历所有使三角形条件满足的 $x$。
\label{thm:biedenharn-elliott}
\end{theorem}
\par\medskip
\noindent\textbf{Biedenharn--Elliott 恒等式的推导}\quad
记 $F$ 为三个角动量之间的再耦合矩阵。由 6j 的定义，$F$ 的每个矩阵元等于一个 6j 符号乘以相位和维数因子。四重张量积从 $(((V_a\otimes V_b)\otimes V_c)\otimes V_d)$ 到 $V_a\otimes(V_b\otimes(V_c\otimes V_d))$ 有两条路径：一条复合三个 $F$ 变换，另一条复合两个 $F$ 变换。
张量积的结合同构必须满足五角形交换图，所以两条复合映射相同。取固定初态、末态的矩阵元，并在三步路径中插入中间不可约表示的完备关系，就得到对允许的 $x$ 求和；每个中间通道贡献其维数 $2x+1$。再代入 $F$ 与 6j 的关系，Condon--Shortley 相位合并为 $(-1)^{R+x}$，正好得到定理中的 Biedenharn--Elliott 恒等式。求和范围无需另行规定：任一三角形条件失败时，相应 6j 符号即为零。
\par\medskip
Biedenharn--Elliott 恒等式就是再耦合的结合律：它保证中间耦合顺序不会改变最终物理矩阵元。
用 Racah 公式可以直接计算具体 6j 符号值。
\medskip
\noindent\emph{$\begin{Bmatrix}\frac12&\frac12&0\\0&0&\frac12\end{Bmatrix}$ 的计算。}
由含零公式（$j_4=0$ 要求 $j_5=j_3=0$, $j_6=j_2=\frac12$）：
\[
\begin{Bmatrix}\tfrac12&\tfrac12&0\\0&0&\tfrac12\end{Bmatrix}
=\frac{(-1)^{j_1+j_2+j_3}}{\sqrt{(2j_2+1)(2j_3+1)}}
=\frac{(-1)^{1/2+1/2+0}}{\sqrt{2\cdot1}}
=-\frac{1}{\sqrt{2}}.
\]
这个负号由 Condon--Shortley 相位约定固定；改变耦合基的整体相位会相应改变单个再耦合矩阵元的符号，因此不能把它单独解释为粒子交换统计。
\par\medskip
\medskip
\noindent\emph{$\begin{Bmatrix}1&\frac12&\frac12\\1&\frac12&\frac12\end{Bmatrix}$ 的逐步计算。}
三角形因子：
\begin{align*}
\Delta(1,\tfrac12,\tfrac12)&=\sqrt{\frac{0!\,1!\,1!}{3!}}=\frac1{\sqrt6},\\
\Delta(1,\tfrac12,\tfrac12)&=\frac1{\sqrt6},\\
\Delta(1,\tfrac12,\tfrac12)&=\frac1{\sqrt6},\\
\Delta(1,\tfrac12,\tfrac12)&=\frac1{\sqrt6}.
\end{align*}
四个 $\Delta$ 之积：$\frac{1}{6\cdot6}=\frac{1}{36}$。
四个三角形和为 $a_1=a_2=a_3=a_4=2$，而三组四角动量和为
$b_1=3,b_2=2,b_3=3$。Racah 求和因而只含 $z=2$ 一项，其值为 $3!=6$；
乘上四个三角形因子得到
\[
\begin{Bmatrix}1&\frac12&\frac12\\1&\frac12&\frac12\end{Bmatrix}
=\frac{6}{36}=\frac16.
\]
四组三角形条件都满足，所以该再耦合矩阵元不应为零。
\par\medskip
\medskip
\noindent\emph{$\begin{Bmatrix}j&j&0\\j&j&0\end{Bmatrix}$ 的一般公式。}
由含零公式：
\[
\begin{Bmatrix}j&j&0\\j&j&0\end{Bmatrix}
=\frac{(-1)^{2j}}{\sqrt{(2j+1)(2j+1)}}
=\frac{(-1)^{2j}}{2j+1}.
\]
对 $j=\frac12$：$\frac{(-1)^1}{2}=-\frac12$。对 $j=1$：$\frac{1}{3}$。对 $j=\frac32$：$\frac{(-1)^3}{4}=-\frac14$。
当两个中间角动量都为零时，数值只剩维数因子 $(2j+1)^{-1}$ 与约定相位 $(-1)^{2j}$。这里的相位来自耦合基约定；半整数角动量并不等同于本式已经实施了一次全同粒子交换。
\par\medskip
当角动量量子数很大时，6j 符号与经典四面体几何有深刻的联系。
\begin{theorem}[Ponzano--Regge 渐近公式]
当所有 $j_i$ 很大时，
\[
\begin{Bmatrix}j_1&j_2&j_3\\j_4&j_5&j_6\end{Bmatrix}
\approx\frac{1}{\sqrt{12\pi V}}
\cos\!\left(\sum_{i=1}^6\left(j_i+\frac12\right)\theta_i+\frac{\pi}{4}\right),
\]
其中 $V$ 是以 $j_1+\frac12,\dots,j_6+\frac12$ 为棱的四面体体积，$\theta_i$ 是棱 $j_i$ 所对二面角。
\label{thm:ponzano-regge}
\end{theorem}
\par\medskip
\noindent\textbf{Ponzano--Regge 公式的推导思路}\quad
四组三角形条件恰好对应四面体的四个面。令
$\ell_i=j_i+\tfrac12$ 为六条棱，并在所有 $j_i$ 以同一尺度趋于无穷时把
Racah 求和作 Euler--Maclaurin 展开。阶乘使用 Stirling 公式
\[
n!\approx\sqrt{2\pi n}\left(\frac{n}{e}\right)^n,\qquad n\to\infty.
\]
代入后得到振荡积分
\[
\begin{Bmatrix}j_1&j_2&j_3\\j_4&j_5&j_6\end{Bmatrix}
\sim\sum_z(-1)^z\exp\bigl[S(z)\bigr]\cdot(\text{prefactor}),
\]
其中 $S(z)$ 是由阶乘对数组成的相位函数。
驻相方程等价于这六条棱构成 Euclidean 四面体。两支驻相点互为复共轭，
驻相相位是 Regge 作用量
$S_{\mathrm R}=\sum_i\ell_i\theta_i$，其中 $\theta_i$ 是外二面角；它不是四面体体积。
二阶涨落的 Hessian 给出振幅 $(12\pi V)^{-1/2}$，两支相加得到
\[
\begin{Bmatrix}j_1&j_2&j_3\\j_4&j_5&j_6\end{Bmatrix}
\sim\frac{1}{\sqrt{12\pi V}}
\cos\left(\sum_i\left(j_i+\frac12\right)\theta_i+\frac{\pi}{4}\right),
\]
其中 $\theta_i$ 是棱 $j_i$ 对面的二面角，$V$ 是四面体体积。这就是 Ponzano--Regge 公式。
这里的相位正是离散三维引力中的单四面体 Regge 作用量。
\par\medskip
Ponzano--Regge 公式把量子角动量再耦合与经典四面体几何联系起来。把 6j 符号视为四面体振幅时，Biedenharn--Elliott 恒等式对应三维流形的 Pachner \(2\text{--}3\) 变换，这也是 Turaev--Viro 不变量和三维离散引力的代数基础。
3j 符号也有 Racah 闭合表达式和超出常规的 Regge 对称性。
\begin{theorem}[3j 符号的 Racah 显式公式]
\[
\begin{pmatrix}j_1&j_2&j_3\\m_1&m_2&m_3\end{pmatrix}
=(-1)^{j_1-j_2-m_3}\Delta(j_1,j_2,j_3)
\sum_z\frac{(-1)^z\sqrt{\prod(j_i\pm m_i)!}}{\text{阶乘组合}(z)},
\]
其中分母与 CG 系数 Racah 公式相同。此公式直接从 \Cref{thm:racah-cg} 和 3j 定义代入得到。
\end{theorem}
\begin{proposition}[Regge 对称性]
3j 符号除了已知的列置换和磁量子数反号外，还满足 Regge 对称性：将 $3j$ 符号编码为 $3\times3$ Regge 阵列
\[
R=\begin{pmatrix}
-j_1+j_2+j_3 & j_1-j_2+j_3 & j_1+j_2-j_3\\
j_1-m_1 & j_2-m_2 & j_3-m_3\\
j_1+m_1 & j_2+m_2 & j_3+m_3
\end{pmatrix},
\]
则 $3j$ 符号在 Regge 阵列的\emph{行置换}和\emph{列置换}下不变（至多差符号 $(-1)^R$，$R$ 为阵列所有元素之和）。这给出 72 种对称操作（$S_3^{\text{row}}\times S_3^{\text{col}}$），远多于之前已知的 12 种。
\end{proposition}
\par\medskip
\noindent\textbf{Regge 对称性的推导}\quad
先Regge 阵列的构造。3j 符号的 Racah 公式为
\[
\begin{pmatrix}j_1&j_2&j_3\\m_1&m_2&m_3\end{pmatrix}
=(-1)^{j_1-j_2-m_3}\Delta(j_1,j_2,j_3)
\sqrt{\prod_{i=1}^{3}(j_i-m_i)!(j_i+m_i)!}\;\sum_z\frac{(-1)^z}{z!\,\prod_k(a_k-z)!\,\prod_l(b_l+z)!},
\]
其中 $a_k$ 取 $j_1+j_2-j_3$、$j_1-m_1$、$j_2+m_2$，$b_l$ 取 $j_2-j_3+m_1$、$j_1-j_3-m_2$、$j_3-j_2-m_1$。分子中九个阶乘的自变量恰好可以排列为 $3\times3$ Regge 阵列 $R$，使每行每列之和都等于 $J=j_1+j_2+j_3$：
\[
R=\begin{pmatrix}
R_{11}&R_{12}&R_{13}\\
R_{21}&R_{22}&R_{23}\\
R_{31}&R_{32}&R_{33}
\end{pmatrix}
=\begin{pmatrix}
-j_1+j_2+j_3 & j_1-j_2+j_3 & j_1+j_2-j_3\\
j_1-m_1 & j_2-m_2 & j_3-m_3\\
j_1+m_1 & j_2+m_2 & j_3+m_3
\end{pmatrix}.
\]
验证行和：第一行 $(-j_1+j_2+j_3)+(j_1-j_2+j_3)+(j_1+j_2-j_3)=j_1+j_2+j_3=J$，同理其余各行。验证列和：第一列 $(-j_1+j_2+j_3)+(j_1-m_1)+(j_1+m_1)=j_1+j_2+j_3=J$，同理其余各列。Racah 公式的分子恰好是 $\sqrt{\prod_{i,j}R_{ij}!}$，分母的三角形因子 $\Delta$ 仅依赖第一行三个元素。
随后行置换不变性。交换 Regge 阵列的两行，对应于 Racah 公式中九个阶乘的重新排列。由于分子中所有阶乘以平方根出现，即 $\sqrt{\prod_{i,j}R_{ij}!}$，行置换只重新排列乘积中的因子，不改变整体乘积值。分母的求和变量 $z$ 的范围由九个阶乘自变量的非负性决定，行置换后求和范围相应变换，但求和值不变（只是求和指标的平移）。唯一可能的变化来自相位 $(-1)^{j_1-j_2-m_3}$ 和 $\Delta$ 因子中第一行元素的变化。对于行置换 $(1\leftrightarrow 2)$，第一行元素 $\{-j_1+j_2+j_3,\,j_1-j_2+j_3,\,j_1+j_2-j_3\}$ 变为 $\{j_1-m_1,\,j_2-m_2,\,j_3-m_3\}$，这对应于 $3j$ 符号中 $(j_i,m_i)$ 参数的某种线性重组合，产生的相位恰好是 $(-1)^R$，其中 $R=\sum_{i,j}R_{ij}=3J$。故行置换至多给出 $(-1)^R$ 相位。
接着列置换不变性。交换 Regge 阵列的两列，对应于 $3j$ 符号的列置换与磁量子数反号的组合。例如交换第 1、2 列：
\[
R_{i1}\leftrightarrow R_{i2}\quad\forall\,i,
\]
即 $-j_1+j_2+j_3\leftrightarrow j_1-j_2+j_3$、$j_1-m_1\leftrightarrow j_2-m_2$、$j_1+m_1\leftrightarrow j_2+m_2$。这恰好对应于 $3j$ 符号中交换 $(j_1,m_1)\leftrightarrow(j_2,m_2)$，即列置换 $(1\leftrightarrow 2)$，已知给出相位 $(-1)^{j_1+j_2-j_3}$。类似地，交换第 1、3 列对应于 $(j_1,m_1)\leftrightarrow(j_3,-m_3)$（注意 $m_3$ 反号），给出已知相位。因此 Regge 阵列的列置换把 $3j$ 符号的列置换与磁量子数反号统一为阵列的列置换操作，使对称性结构更清晰。
最后，$3\times3$ 阵列的行、列置换给出 $6\times6=36$ 个操作；转置 $R\mapsto R^{\mathsf T}$ 再给出另一陪集。因此总对称数为 $72$。
\par\medskip
这些操作组成半直积 $(S_3^{\mathrm{row}}\times S_3^{\mathrm{col}})\rtimes\ZZ_2$；转置生成元交换行、列两个 $S_3$ 因子，所以不能写成三个群的直积。其阶数为
\begin{equation*}
|G_{\text{Regge}}|=|S_3^{\text{row}}|\times|S_3^{\text{col}}|\times|\ZZ_2^{\text{transpose}}|=6\times6\times2=72.
\end{equation*}
它同构于 wreath 积 $S_3\wr\ZZ_2$。在数值计算中，可利用这一轨道结构把等价的 3j 符号统一映射到同一标准代表，并把对称性作为结果校验。
\section{不可约张量算符与 Wigner--Eckart 定理}
\label{sec:wigner-eckart}
前面各节解决了态空间的角动量耦合问题：如何把 $j_1\otimes j_2$ 分解为不可约表示的直和。另一个同样基本的问题是\textbf{算符的矩阵元}。在量子力学中，一切可观测量（能级、跃迁概率、散射振幅）都归结为算符在态之间的矩阵元 $\bra{j',m'}T\ket{j,m}$。当系统具有转动对称性时，这些矩阵元不能任意取值；它们被群论严格约束，可以分解为一个纯几何因子（$3j$ 符号）和一个动力学因子（约化矩阵元）。这就是 Wigner--Eckart 定理，它是表示论在量子力学中最有力的工具之一。
\begin{definitionplain}[不可约张量算符]
一组 $2k+1$ 个算符 $\{T^{(k)}_q\mid q=-k,-k+1,\dots,k\}$ 称为\textbf{ 秩 $k$ 的不可约张量算符}（或球张量算符），如果在转动 $U(R)$ 下它们按 $SO(3)$ 的不可约表示 $D^{(k)}$ 变换：
\begin{equation}
U(R)\,T^{(k)}_q\,U(R)^\dagger
=\sum_{q'=-k}^{k}D^{(k)}_{q'q}(R)\,T^{(k)}_{q'}.
\label{eq:tensor-op-def}
\end{equation}
\end{definitionplain}
这个定义的物理含义是：$T^{(k)}_q$ 作为一个算符集合，在转动下像秩-$k$ 的角动量态一样变换。对比态的变换规律 $U(R)\ket{j,m}=\sum_{m'}D^{(j)}_{m'm}(R)\ket{j,m'}$，可见算符的变换是\textbf{共变}的（$U$ 在左，$U^\dagger$ 在右），而态的变换是\textbf{逆变}的。这种结构保证了矩阵元 $\bra{j',m'}T^{(k)}_q\ket{j,m}$ 在转动下不变——它是三个表示 $D^{(j')*}\otimes D^{(k)}\otimes D^{(j)}$ 中的标量分量。
\begin{proposition}[无穷小形式：对易关系判据]
算符集合 $\{T^{(k)}_q\}$ 是秩 $k$ 不可约张量算符，当且仅当它与角动量算符满足对易关系
\begin{align}
[J_z,T^{(k)}_q]&=\hbar q\,T^{(k)}_q,\label{eq:tensor-Jz}\\
[J_\pm,T^{(k)}_q]&=\hbar\sqrt{k(k+1)-q(q\pm1)}\,T^{(k)}_{q\pm1}.\label{eq:tensor-Jpm}
\end{align}
\end{proposition}
\begin{proof}
取 $R=R_z(\theta)$（绕 $z$ 轴转角 $\theta$），则 $U(R)=\ee^{-\ii J_z\theta/\hbar}$，代入 \cref{eq:tensor-op-def}：
\[
\ee^{-\ii J_z\theta/\hbar}T^{(k)}_q\ee^{\ii J_z\theta/\hbar}
=\sum_{q'}D^{(k)}_{q'q}(R_z(\theta))T^{(k)}_{q'}
=\ee^{-\ii q\theta}T^{(k)}_q,
\]
因为 $D^{(k)}$ 在 $R_z$ 下是对角的：$D^{(k)}_{q'q}(R_z(\theta))=\delta_{q'q}\ee^{-\ii q\theta}$。左边展开到 $\theta$ 一阶：
\[
T^{(k)}_q-\frac{\ii\theta}{\hbar}[J_z,T^{(k)}_q]+O(\theta^2)
=T^{(k)}_q-\ii q\theta\,T^{(k)}_q+O(\theta^2),
\]
比较 $\theta$ 一阶项得 \cref{eq:tensor-Jz}。类似地，取无穷小转动绕 $x$ 和 $y$ 轴，组合成 $J_\pm=J_x\pm\ii J_y$，利用 $D^{(k)}$ 矩阵元
\[
D^{(k)}_{q',q}(R_\pm(\delta\theta))=\delta_{q',q\pm1}\sqrt{k(k+1)-q(q\pm1)}\,\delta\theta+O(\delta\theta^2),
\]
得到 \cref{eq:tensor-Jpm}。逆方向：从对易关系出发，用 Baker--Campbell--Hausdorff 公式重建有限转动下的变换规律 \cref{eq:tensor-op-def}。
\end{proof}
对易关系 \cref{eq:tensor-Jz} 有一个直接的物理解读：$T^{(k)}_q$ 将态的磁量子数改变 $q$。因为若 $\ket{j,m}$ 是 $J_z$ 的本征态，则
\[
J_z\bigl(T^{(k)}_q\ket{j,m}\bigr)
=\bigl([J_z,T^{(k)}_q]+T^{(k)}_qJ_z\bigr)\ket{j,m}
=\hbar(q+m)T^{(k)}_q\ket{j,m},
\]
所以 $T^{(k)}_q\ket{j,m}$ 是 $J_z$ 本征值为 $\hbar(m+q)$ 的态（若非零）。这是第一条选择定则的来源。
\medskip
\noindent\emph{标量算符（秩 $0$）。}
秩 $0$ 的不可约张量算符只有一个分量 $T^{(0)}_0$，对易关系给出 $[J_i,T^{(0)}_0]=0$（与所有角动量分量对易）。因此秩 $0$ 张量算符就是转动不变量（标量），例如 $\vb J^2$、$\vb r^2$、$\vb p^2$、中心势 $V(r)$、Hamilton 量 $\hat H$（对转动不变系统）。
\par\medskip
\medskip
\noindent\emph{矢量算符（秩 $1$）。}
秩 $1$ 的不可约张量算符有三个分量 $T^{(1)}_{+1},T^{(1)}_0,T^{(1)}_{-1}$，与矢量 $\vb V=(V_x,V_y,V_z)$ 的关系为
\[
T^{(1)}_{\pm1}=\mp\frac{1}{\sqrt2}(V_x\pm\ii V_y),\qquad T^{(1)}_0=V_z.
\]
这恰好是球基矢 $\vb e_{\pm1}=\mp(\vb e_x\pm\ii\vb e_y)/\sqrt2$、$\vb e_0=\vb e_z$ 下的分量：$T^{(1)}_q=\vb V\cdot\vb e_q^*$。典型矢量算符包括位置 $\vb r$、动量 $\vb p$、角动量 $\vb J$、自旋 $\vb S$、电偶极矩 $\vb d=-e\vb r$。
验证对易关系：对 $\vb r$，$[J_z,z]=0$（因 $J_z=xp_y-yp_x$ 与 $z$ 对易），$[J_z,x\pm\ii y]=\pm\ii\hbar(x\pm\ii y)$，故 $[J_z,r^{(1)}_q]=\hbar q\,r^{(1)}_q$。类似地 $[J_\pm,z]=\pm\hbar(x\pm\ii y)/\sqrt2\cdot\sqrt2=\hbar(\mp\sqrt2\,r^{(1)}_{\pm1})$，与 \cref{eq:tensor-Jpm}（$k=1$）一致。
\par\medskip
\medskip
\noindent\emph{球谐函数作为张量算符。}
球谐函数 $Y_{kq}(\theta,\varphi)$ 作为乘法算符（作用在波函数上：$\psi\mapsto Y_{kq}\psi$）构成秩 $k$ 不可约张量算符。这是因为球谐函数在转动下按 $D^{(k)}$ 变换：
\[
Y_{kq}(R^{-1}\uv r)=\sum_{q'}D^{(k)}_{q'q}(R)Y_{kq'}(\uv r),
\]
而主动转动下波函数变换为 $\psi'(r)=\psi(R^{-1}r)$，故 $Y_{kq}$ 作为算符的变换规律恰好是 \cref{eq:tensor-op-def}。特别地：
\begin{itemize}
\item $Y_{00}=1/\sqrt{4\pi}$：标量（秩 $0$）。
\item $Y_{1q}\propto r^{(1)}_q/r$：归一化的位置分量（秩 $1$）。
\item $Y_{2q}$：四极矩算符（秩 $2$），如 $Y_{20}\propto(3z^2-r^2)/r^2$、$Y_{2,\pm2}\propto(x\pm\ii y)^2/r^2$。
\end{itemize}
\par\medskip
\medskip
\noindent\emph{四极矩算符（秩 $2$）。}
电四极矩张量 $Q_{ij}=e(3r_ir_j-r^2\delta_{ij})$ 是对称无迹笛卡儿张量，有 5 个独立分量。对应的球张量分量为
\[
Q^{(2)}_0=2ez^2-e(x^2+y^2)=er^2(3\cos^2\theta-1),
\]
\[
Q^{(2)}_{\pm1}=\mp\frac{\sqrt6}{2}e z(x\pm\ii y),\qquad
Q^{(2)}_{\pm2}=\frac{\sqrt6}{4}e(x\pm\ii y)^2.
\]
这 5 个分量在转动下按 $D^{(2)}$ 变换，构成秩 $2$ 不可约张量算符。在核物理中，原子核的四极矩 $Q$ 定义为
\[
Q=\sqrt{\frac{16\pi}{5}}\braket{I,I}{Q^{(2)}_0}{I,I},
\]
其中 $I$ 为核自旋，它直接决定超精细结构中的电四极相互作用强度。
\par\medskip
\begin{theorem}[Wigner--Eckart 定理]
\label{thm:wigner-eckart-angular}
设 $T^{(k)}_q$ 是秩 $k$ 不可约张量算符，$\ket{j,m}$ 是角动量本征态。则矩阵元
\begin{equation}
\bra{j',m'}T^{(k)}_q\ket{j,m}
=(-1)^{j'-m'}\begin{pmatrix}j'&k&j\\-m'&q&m\end{pmatrix}\braket{j'\|T^{(k)}\|j},
\label{eq:wigner-eckart}
\end{equation}
其中 $\braket{j'\|T^{(k)}\|j}$ 是\textbf{约化矩阵元}（约化矩阵单元），只依赖 $j,j',k$ 和算符的物理内容，不依赖磁量子数 $m,m',q$。
\end{theorem}
这个定理的含义是：矩阵元被分解为两个因子的乘积——$3j$ 符号是纯几何因子（完全由角动量耦合规则决定），约化矩阵元是纯动力学因子（包含物理实质）。所有 $m,m',q$ 的依赖都被 $3j$ 符号完全确定。对于固定的 $j,j',k$，共有 $(2j+1)(2j'+1)(2k+1)$ 个矩阵元，但它们全部由一个数 $\braket{j'\|T^{(k)}\|j}$ 确定。
\par\medskip
\noindent\textbf{Wigner--Eckart 定理的证明}\quad
先构造不变量。考虑态 $T^{(k)}_q\ket{j,m}$。这个态在 $J_z$ 下的本征值为 $\hbar(m+q)$（由对易关系 \cref{eq:tensor-Jz}），在 $J^2$ 下不一定有确定值，但可以按总角动量本征态展开：
\[
T^{(k)}_q\ket{j,m}=\sum_{J,M}c_{J,M}(q,m)\ket{J,M},
\]
其中 $M=m+q$（$J_z$ 本征值约束）。
随后利用转动不变性确定展开系数。对任意转动 $R$，利用 \cref{eq:tensor-op-def}：
\[
U(R)T^{(k)}_q\ket{j,m}
=\sum_{q',m'}D^{(k)}_{q'q}(R)D^{(j)}_{m'm}(R)T^{(k)}_{q'}\ket{j,m'}.
\]
另一方面，
\[
U(R)\sum_{J,M}c_{J,M}(q,m)\ket{J,M}
=\sum_{J,M,M'}c_{J,M}(q,m)D^{(J)}_{M'M}(R)\ket{J,M'}.
\]
比较两边：左边的 $D^{(k)}\otimes D^{(j)}$ 表示在右边被分解为 $D^{(J)}$ 的直和。由 Clebsch--Gordan 分解 \cref{eq:cg-decomposition}，展开系数必须正比于 CG 系数：
\[
c_{J,M}(q,m)=\alpha(J)\,C(k,j,J;q,m,M),
\]
其中 $\alpha(J)$ 是比例常数，不依赖 $q,m,M$（因为 Schur 引理保证在同一不可约表示内比例因子是标量）。因此
\begin{equation}
T^{(k)}_q\ket{j,m}=\sum_{J=|k-j|}^{k+j}\alpha(J)\sum_{M}C(k,j,J;q,m,M)\ket{J,M}.
\label{eq:tensor-op-expansion}
\end{equation}
接着提取矩阵元。取 $\bra{j',m'}$ 内积：
\[
\bra{j',m'}T^{(k)}_q\ket{j,m}
=\sum_J\alpha(J)\sum_M C(k,j,J;q,m,M)\braket{j',m'}{J,M}
=\alpha(j')\,C(k,j,j';q,m,m'),
\]
利用 $\braket{j',m'}{J,M}=\delta_{j'J}\delta_{m'M}$。CG 系数与 $3j$ 符号的关系为
\[
C(k,j,j';q,m,m')=(-1)^{k-j+m'}\sqrt{2j'+1}\begin{pmatrix}k&j&j'\\q&m&m'\end{pmatrix}.
\]
利用 $3j$ 符号的列置换对称性（交换列 1 和 3 给出 $(-1)^{k+j+j'}$）：
\[
\begin{pmatrix}k&j&j'\\q&m&m'\end{pmatrix}
=(-1)^{k+j+j'}\begin{pmatrix}j'&k&j\\m'&q&m\end{pmatrix}
=(-1)^{k+j+j'}(-1)^{j'-m'}\begin{pmatrix}j'&k&j\\-m'&q&m\end{pmatrix},
\]
标准的 CG--$3j$ 约定为：
\[
\bra{j',m'}T^{(k)}_q\ket{j,m}
=\frac{\alpha(j')}{\sqrt{2j'+1}}(-1)^{j'-m'}\begin{pmatrix}j'&k&j\\-m'&q&m\end{pmatrix}\cdot(2j'+1),
\]
整理后定义约化矩阵元
\[
\braket{j'\|T^{(k)}\|j}:=\sqrt{2j'+1}\,\alpha(j'),
\]
即得 \cref{eq:wigner-eckart}。
还需验证约化矩阵元不依赖磁量子数。取两组满足
\[
m'_1=m_1+q_1,
\qquad
m'_2=m_2+q_2
\]
的磁量子数，并分别写出相应矩阵元。由 \cref{eq:wigner-eckart}：
\[
\frac{\bra{j',m'_1}T^{(k)}_{q_1}\ket{j,m_1}}{(-1)^{j'-m'_1}\begin{pmatrix}j'&k&j\\-m'_1&q_1&m_1\end{pmatrix}}
=
\frac{\bra{j',m'_2}T^{(k)}_{q_2}\ket{j,m_2}}{(-1)^{j'-m'_2}\begin{pmatrix}j'&k&j\\-m'_2&q_2&m_2\end{pmatrix}}
=\braket{j'\|T^{(k)}\|j},
\]
即任意一个非零矩阵元除以对应的几何因子都给出同一个约化矩阵元。这在实践中极为有用：只需计算一个最方便的矩阵元，就能确定所有其他矩阵元。
\par\medskip
Wigner--Eckart 定理立即给出两条基本选择定则：
\begin{proposition}[角动量选择定则]
矩阵元 $\bra{j',m'}T^{(k)}_q\ket{j,m}$ 非零的必要条件是：
\begin{enumerate}
\item \textbf{三角形条件}：$|j-k|\le j'\le j+k$（等价于 $j,j',k$ 构成三角形三边）。
\item \textbf{磁量子数条件}：$m'=m+q$。
\end{enumerate}
\end{proposition}
\begin{proof}
$3j$ 符号 $\begin{pmatrix}j'&k&j\\-m'&q&m\end{pmatrix}$ 非零要求三个 $j$ 满足三角形条件（来自 CG 系数的定义域约束）且三个 $m$ 之和为零：$-m'+q+m=0$，即 $m'=m+q$。
\end{proof}
\begin{proposition}[宇称选择定则]
设系统具有空间反演对称性，$\ket{j,m}$ 的宇称为 $\pi=(-1)^l$（$l$ 为轨道角动量），秩 $k$ 张量算符 $T^{(k)}$ 的宇称为 $\pi_T$（电多极 $\pi_T=(-1)^k$，磁多极 $\pi_T=(-1)^{k+1}$）。则矩阵元非零要求
\[
\pi'\pi_T\pi=+1,
\]
即初态、算符、末态宇称之积为偶。这就是 Laporte 定则的一般形式。
\end{proposition}
\begin{proof}
反演算符 $P$ 与转动对易，故 $P\ket{j,m}=\pi\ket{j,m}$，$PT^{(k)}_qP^{-1}=\pi_T T^{(k)}_q$。矩阵元在反演下不变：
\[
\bra{j',m'}T^{(k)}_q\ket{j,m}
=\bra{j',m'}P^{-1}PT^{(k)}_qP^{-1}P\ket{j,m}
=\pi'\pi_T\pi\,\bra{j',m'}T^{(k)}_q\ket{j,m},
\]
故 $(1-\pi'\pi_T\pi)\bra{j',m'}T^{(k)}_q\ket{j,m}=0$。若 $\pi'\pi_T\pi=-1$ 则矩阵元为零。
\end{proof}
\medskip
\noindent\emph{电偶极跃迁的选择定则。}
电偶极矩 $\vb d=-e\vb r$ 是秩 $1$ 张量算符（$k=1$），宇称 $\pi_T=-1$（奇宇称，因 $\vb r\to-\vb r$）。选择定则：
\begin{itemize}
\item $\Delta l=\pm1$（三角形条件 $|l-1|\le l'\le l+1$ 加上宇称 $(-1)^{l'}(-1)(-1)^l=+1$ 要求 $l'+l$ 为奇数，排除 $l'=l$）。
\item $\Delta m=0,\pm1$（磁量子数条件 $m'=m+q$，$q=0,\pm1$）。
\item $\Delta j=0,\pm1$（但 $j=0\to j'=0$ 禁戒，因 $3j$ 符号 $\begin{pmatrix}0&1&0\\0&q&0\end{pmatrix}=0$）。
\end{itemize}
这些正是原子光谱中熟知的电偶极跃迁选择定则。注意 $\Delta l=\pm1$ 来自宇称，而 $\Delta j=0,\pm1$ 来自角动量三角形条件——两者是独立的约束。
\par\medskip
Wigner--Eckart 定理把所有矩阵元的计算归结为一个约化矩阵元。实践中，约化矩阵元可以通过以下方法确定：
\begin{proposition}[约化矩阵元的确定方法]
若存在一组 $(m_0,q_0,m'_0)$ 使得 $3j$ 符号 $\begin{pmatrix}j'&k&j\\-m'_0&q_0&m_0\end{pmatrix}\ne0$，则
\begin{equation}
\braket{j'\|T^{(k)}\|j}
=\frac{\bra{j',m'_0}T^{(k)}_{q_0}\ket{j,m_0}}{(-1)^{j'-m'_0}\begin{pmatrix}j'&k&j\\-m'_0&q_0&m_0\end{pmatrix}}.
\label{eq:reduced-me-formula}
\end{equation}
通常选取 $m_0=j$，$q_0=m'_0-j$（使计算最简单的组合）。
\end{proposition}
\medskip
\noindent\emph{位置算符在 $l=1$ 子空间的约化矩阵元。}
计算 $\braket{l'=1\|r^{(1)}\|l=1}$。取 $m_0=1$，$q_0=0$，$m'_0=1$：
\[
\bra{1,1}r^{(1)}_0\ket{1,1}=\bra{1,1}z\ket{1,1}=\int_0^\infty r^3\,dr\int|Y_{11}|^2\cos\theta\,d\Omega.
\]
利用 $Y_{11}=-\sqrt{3/8\pi}\sin\theta\,\ee^{\ii\varphi}$ 和 $\cos\theta=\sqrt{4\pi/3}Y_{10}$：
\[
\int Y_{11}^*\cos\theta\,Y_{11}\,d\Omega
=\sqrt{\frac{4\pi}{3}}\int Y_{11}^*Y_{10}Y_{11}\,d\Omega
=\sqrt{\frac{4\pi}{3}}\cdot\left(-\sqrt{\frac{3}{8\pi}}\right)\sqrt{\frac{3}{8\pi}}\sqrt{\frac{3}{4\pi}}\int\sin^2\theta\cos\theta\sin\theta\,d\theta\,d\varphi.
\]
更简洁地，利用 Gaunt 积分公式 $\int Y_{l_1m_1}Y_{l_2m_2}Y_{l_3m_3}\,d\Omega$ 正比于 $3j$ 符号 $\begin{pmatrix}l_1&l_2&l_3\\m_1&m_2&m_3\end{pmatrix}$，最终得到
\[
\braket{1\|r^{(1)}\|1}=\frac{\bra{1,1}z\ket{1,1}}{(-1)^0\begin{pmatrix}1&1&1\\-1&0&1\end{pmatrix}}=r\sqrt{\frac{3}{2}},
\]
其中径向积分 $\int_0^\infty r^3|R_{nl}(r)|^2\,dr$ 已简记为 $r$。由此可立即写出所有 $\bra{1,m'}r^{(1)}_q\ket{1,m}$。
\par\medskip
\begin{theorem}[投影定理（等价算符方法）]
\label{thm:projection}
设 $\vb V$ 是矢量算符（秩 $1$ 张量），在角动量 $j$ 的不可约子空间内（$j\ne0$），$\vb V$ 的矩阵元与 $\vb J$ 的矩阵元成正比：
\begin{equation}
\bra{j,m'}\vb V\ket{j,m}
=\frac{\braket{\vb V\cdot\vb J}}{\hbar^2 j(j+1)}\bra{j,m'}\vb J\ket{j,m},
\label{eq:projection-theorem}
\end{equation}
其中 $\braket{\vb V\cdot\vb J}:=\bra{j,m}\vb V\cdot\vb J\ket{j,m}$（不依赖 $m$，因 $\vb V\cdot\vb J$ 是标量）。
\end{theorem}
\begin{proof}
$\vb V\cdot\vb J$ 是秩-$0$ 标量，故由 Wigner--Eckart 定理，其对角矩阵元不依赖 $m$：
\[
\braket{j,m}{\vb V\cdot\vb J}{j,m}=\braket{j,m'}{\vb V\cdot\vb J}{j,m'}\quad\forall\,m,m'.
\]
记此公共值为 $\braket{\vb V\cdot\vb J}$。
在同一 $j$ 子空间内，$\vb V$ 和 $\vb J$ 又都是秩-$1$ 张量算符，所以矩阵元正比于同一个 $3j$ 符号：
\[
\bra{j,m'}V^{(1)}_q\ket{j,m}=(-1)^{j-m'}\begin{pmatrix}j&1&j\\-m'&q&m\end{pmatrix}\braket{j\|V^{(1)}\|j},
\]
\[
\bra{j,m'}J^{(1)}_q\ket{j,m}=(-1)^{j-m'}\begin{pmatrix}j&1&j\\-m'&q&m\end{pmatrix}\braket{j\|J^{(1)}\|j}.
\]
因此 $\bra{j,m'}V^{(1)}_q\ket{j,m}=\frac{\braket{j\|V^{(1)}\|j}}{\braket{j\|J^{(1)}\|j}}\bra{j,m'}J^{(1)}_q\ket{j,m}$。
比例常数由标量积 $\vb V\cdot\vb J=\sum_q(-1)^qV^{(1)}_qJ^{(1)}_{-q}$ 的对角矩阵元确定：
\[
\braket{\vb V\cdot\vb J}=\sum_q(-1)^q\bra{j,m}V^{(1)}_qJ^{(1)}_{-q}\ket{j,m}.
\]
插入完备集 $\sum_{m'}\ket{j,m'}\bra{j,m'}$（同一 $j$ 子空间内）：
\[
\braket{\vb V\cdot\vb J}=\sum_{m',q}(-1)^q\bra{j,m}V^{(1)}_q\ket{j,m'}\bra{j,m'}J^{(1)}_{-q}\ket{j,m}.
\]
用 Wigner--Eckart 定理代入两个矩阵元。$3j$ 符号的正交关系为
\[
\sum_{m_1,m_2}
\begin{pmatrix}j_1&j_2&j_3\\m_1&m_2&m_3\end{pmatrix}
\begin{pmatrix}j_1&j_2&j_3'\\m_1&m_2&m_3'\end{pmatrix}
=\frac{\delta_{j_3j_3'}\delta_{m_3m_3'}}{2j_3+1};
\]
此处 $j_2=1$，正交化因子为 $1/(2\cdot1+1)=1/3$，得到
\[
\braket{\vb V\cdot\vb J}=\frac{1}{3}\braket{j\|V^{(1)}\|j}\braket{j\|J^{(1)}\|j}.
\]
类似地 $\braket{\vb J\cdot\vb J}=\hbar^2 j(j+1)=\frac{1}{3}\braket{j\|J^{(1)}\|j}^2$，故 $\braket{j\|J^{(1)}\|j}=\hbar\sqrt{3j(j+1)}$，代入得比例常数
\[
\frac{\braket{j\|V^{(1)}\|j}}{\braket{j\|J^{(1)}\|j}}=\frac{\braket{\vb V\cdot\vb J}}{\hbar^2 j(j+1)}.
\]
\end{proof}
\medskip
\noindent\emph{Land\'e $g$ 因子。}
电子的总磁矩 $\vb\mu=-\mu_B(g_L\vb L+g_S\vb S)/\hbar$，其中 $g_L=1$（轨道），$g_S\approx2$（自旋）。在总角动量 $\vb J=\vb L+\vb S$ 的本征态 $\ket{j,m_j}$ 中，磁矩的期望值需要计算 $\bra{j,m_j}\vb\mu\ket{j,m_j}$。
由投影定理 \cref{thm:projection}，在同一 $j$ 子空间内：
\[
\vb\mu\to\frac{\braket{\vb\mu\cdot\vb J}}{\hbar^2 j(j+1)}\vb J.
\]
计算 $\vb\mu\cdot\vb J$：
\[
\vb\mu\cdot\vb J=-\frac{\mu_B}{\hbar}(g_L\vb L\cdot\vb J+g_S\vb S\cdot\vb J).
\]
利用 $\vb L\cdot\vb J=\frac12(\vb J^2+\vb L^2-\vb S^2)$ 和 $\vb S\cdot\vb J=\frac12(\vb J^2+\vb S^2-\vb L^2)$，在 $\ket{l,s,j,m_j}$ 态中：
\[
\braket{\vb L\cdot\vb J}=\frac{\hbar^2}{2}[j(j+1)+l(l+1)-s(s+1)],
\]
\[
\braket{\vb S\cdot\vb J}=\frac{\hbar^2}{2}[j(j+1)+s(s+1)-l(l+1)].
\]
代入得 Land\'e $g$ 因子：
\begin{equation}
g_J=1+\frac{j(j+1)+s(s+1)-l(l+1)}{2j(j+1)},
\label{eq:lande-g}
\end{equation}
Zeeman 效应中的能量修正为 $\Delta E=\mu_B g_J m_j B$。这个公式是投影定理最经典的应用——群论把一个需要计算大量矩阵元的问题简化为一个代数表达式。
\par\medskip
两个不可约张量算符可以耦合为新的张量算符，正如两个角动量态可以耦合为总角动量态。
\begin{proposition}[张量算符的耦合]
设 $T^{(k_1)}_{q_1}$ 和 $S^{(k_2)}_{q_2}$ 分别是秩 $k_1$ 和 $k_2$ 的不可约张量算符。定义耦合张量算符
\begin{equation}
[T^{(k_1)}\otimes S^{(k_2)}]^{(K)}_Q
=\sum_{q_1,q_2}C(k_1,k_2,K;q_1,q_2,Q)\,T^{(k_1)}_{q_1}S^{(k_2)}_{q_2},
\label{eq:tensor-product-op}
\end{equation}
则 $\{[T^{(k_1)}\otimes S^{(k_2)}]^{(K)}_Q\}_{Q=-K}^K$ 是秩 $K$ 不可约张量算符，其中 $K$ 取值 $|k_1-k_2|,\dots,k_1+k_2$。
\end{proposition}
\begin{proof}
需要验证耦合算符满足对易关系 \cref{eq:tensor-Jz} 和 \cref{eq:tensor-Jpm}。利用 Leibniz 规则 $[J_z,AB]=[J_z,A]B+A[J_z,B]$：
\[
[J_z,T^{(k_1)}_{q_1}S^{(k_2)}_{q_2}]
=\hbar(q_1+q_2)T^{(k_1)}_{q_1}S^{(k_2)}_{q_2}.
\]
CG 系数 $C(k_1,k_2,K;q_1,q_2,Q)$ 非零要求 $q_1+q_2=Q$，故
\[
[J_z,[T^{(k_1)}\otimes S^{(k_2)}]^{(K)}_Q]
=\hbar Q\,[T^{(k_1)}\otimes S^{(k_2)}]^{(K)}_Q.
\]
类似地，$J_\pm$ 的对易关系利用 CG 系数的递推关系
\[
\sum_{q_1,q_2}C(k_1,k_2,K;q_1,q_2,Q)\bigl[\sqrt{k_1(k_1+1)-q_1(q_1\pm1)}\,\delta_{q_1'\,q_1\pm1}+\sqrt{k_2(k_2+1)-q_2(q_2\pm1)}\,\delta_{q_2'\,q_2\pm1}\bigr]
\]
\[
=\sqrt{K(K+1)-Q(Q\pm1)}\,C(k_1,k_2,K;q_1',q_2',Q\pm1),
\]
这正是 CG 系数作为角动量耦合系数所满足的递推关系。故 $[T^{(k_1)}\otimes S^{(k_2)}]^{(K)}_Q$ 满足 \cref{eq:tensor-Jpm}。
\end{proof}
\begin{proposition}[标量积的矩阵元]
两个同秩张量算符的标量积定义为
\[
T^{(k)}\cdot S^{(k)}:=\sum_{q=-k}^{k}(-1)^q T^{(k)}_q S^{(k)}_{-q}
=(-1)^k\sqrt{2k+1}\,[T^{(k)}\otimes S^{(k)}]^{(0)}_0.
\]
其矩阵元为
\begin{equation}
\bra{j',m'}T^{(k)}\cdot S^{(k)}\ket{j,m}
=(-1)^{j'-j}\delta_{m'm}\frac{\braket{j'\|T^{(k)}\|j}\braket{j'\|S^{(k)}\|j}}{\sqrt{2j'+1}}\,\delta_{j'j},
\label{eq:scalar-product-me}
\end{equation}
当 $j'=j$ 时（标量积不改变角动量），进一步简化为
\[
\bra{j,m}T^{(k)}\cdot S^{(k)}\ket{j,m}
=\frac{1}{2j+1}\braket{j\|T^{(k)}\|j}\braket{j\|S^{(k)}\|j}.
\]
\end{proposition}
\begin{proof}
标量积是秩 $0$ 张量，由 Wigner--Eckart 定理：
\[
\bra{j',m'}[T^{(k)}\otimes S^{(k)}]^{(0)}_0\ket{j,m}
=(-1)^{j'-m'}\begin{pmatrix}j'&0&j\\-m'&0&m\end{pmatrix}\braket{j'\|[T^{(k)}\otimes S^{(k)}]^{(0)}\|j}.
\]
$3j$ 符号 $\begin{pmatrix}j'&0&j\\-m'&0&m\end{pmatrix}=\frac{(-1)^{j'-m'}}{\sqrt{2j'+1}}\delta_{j'j}\delta_{m'm}$（秩 $0$ 耦合是平凡的）。约化矩阵元的分解涉及两个张量算符约化矩阵元的乘积和一个 $6j$ 符号：
\[
\braket{j'\|[T^{(k)}\otimes S^{(k)}]^{(0)}\|j}
=(-1)^{j+k+j'}\sqrt{2j'+1}\begin{Bmatrix}j&k&j'\\k&j&0\end{Bmatrix}\braket{j'\|T^{(k)}\|j}\braket{j'\|S^{(k)}\|j}.
\]
含 $0$ 的 $6j$ 符号简化为 $\begin{Bmatrix}j&k&j'\\k&j&0\end{Bmatrix}=\frac{(-1)^{j+k+j'}}{\sqrt{(2j+1)(2k+1)}}\delta_{jj'}$（见 \cref{eq:6j-zero}），代入整理即得 \cref{eq:scalar-product-me}。
\end{proof}
\medskip
\noindent\emph{自旋--轨道耦合的矩阵元。}
自旋--轨道耦合 $\vb L\cdot\vb S$ 是两个秩 $1$ 张量算符的标量积。在 $\ket{l,s,j,m_j}$ 态中：
\[
\vb L\cdot\vb S=\frac12(\vb J^2-\vb L^2-\vb S^2),
\]
故
\[
\bra{l,s,j,m_j}\vb L\cdot\vb S\ket{l,s,j,m_j}
=\frac{\hbar^2}{2}[j(j+1)-l(l+1)-s(s+1)].
\]
用标量积公式 \cref{eq:scalar-product-me} 验证：需要 $\braket{j\|L^{(1)}\|j}$ 和 $\braket{j\|S^{(1)}\|j}$。由投影定理，$\vb L\to\frac{\braket{\vb L\cdot\vb J}}{\hbar^2 j(j+1)}\vb J$，$\vb S\to\frac{\braket{\vb S\cdot\vb J}}{\hbar^2 j(j+1)}\vb J$，故
\[
\braket{j\|L^{(1)}\|j}=\frac{\braket{\vb L\cdot\vb J}}{\hbar^2 j(j+1)}\braket{j\|J^{(1)}\|j},
\]
对 $\vb S$ 同理。标量积公式给出
\[
\begin{aligned}
\bra{j,m_j}\vb L\cdot\vb S\ket{j,m_j}
&=\frac{\braket{\vb L\cdot\vb J}\braket{\vb S\cdot\vb J}}
{(2j+1)\hbar^4j^2(j+1)^2}\,
\braket{j\|J^{(1)}\|j}^2\\
&=\frac{\braket{\vb L\cdot\vb J}\braket{\vb S\cdot\vb J}}
{\hbar^2j(j+1)}.
\end{aligned}
\]
与直接计算 $\frac12[j(j+1)-l(l+1)-s(s+1)]\hbar^2$ 一致（利用 $\braket{\vb L\cdot\vb J}+\braket{\vb S\cdot\vb J}=\hbar^2 j(j+1)$）。
\par\medskip
当三个张量算符依次作用于态上时，不同耦合顺序之间的变换系数正是 $6j$ 符号。这是 $6j$ 符号在物理中最重要的出现方式之一。
\begin{theorem}[三个张量算符的约化矩阵元公式]
设 $T^{(k_1)}$ 和 $S^{(k_2)}$ 是不可约张量算符。耦合张量 $[T^{(k_1)}\otimes S^{(k_2)}]^{(K)}$ 的约化矩阵元为
\begin{equation}
\begin{aligned}
&\braket{j'\|[T^{(k_1)}\otimes S^{(k_2)}]^{(K)}\|j}\\
&\quad=\sum_{j''}(-1)^{j+j'+j''}
\sqrt{(2K+1)(2j''+1)}
\begin{Bmatrix}k_1&k_2&K\\j&j'&j''\end{Bmatrix}\\
&\qquad\times
\braket{j'\|T^{(k_1)}\|j''}
\braket{j''\|S^{(k_2)}\|j}.
\end{aligned}
\label{eq:reduced-me-product}
\end{equation}
\end{theorem}
\begin{proof}
思路。将 $[T^{(k_1)}\otimes S^{(k_2)}]^{(K)}_Q\ket{j,m}$ 按 \cref{eq:tensor-op-expansion} 展开，中间态 $\ket{j'',m''}$ 来自 $S^{(k_2)}$ 作用于 $\ket{j,m}$，然后 $T^{(k_1)}$ 作用于 $\ket{j'',m''}$。两次耦合的顺序 $(k_2\text{ 先}, k_1\text{ 后})$ 与直接耦合 $(k_1\otimes k_2\to K)$ 之间的变换系数正是 $6j$ 符号。
详细推导。由 Wigner--Eckart 定理，
\[
\bra{j',m'}[T^{(k_1)}\otimes S^{(k_2)}]^{(K)}_Q\ket{j,m}
=(-1)^{j'-m'}\begin{pmatrix}j'&K&j\\-m'&Q&m\end{pmatrix}\braket{j'\|[T^{(k_1)}\otimes S^{(k_2)}]^{(K)}\|j}.
\]
另一方面，直接展开耦合算符 \cref{eq:tensor-product-op} 并用 Wigner--Eckart 定理分别处理两个矩阵元：
\begin{align*}
&\bra{j',m'}[T^{(k_1)}\otimes S^{(k_2)}]^{(K)}_Q\ket{j,m}\\
&\quad=\sum_{q_1,q_2}C(k_1,k_2,K;q_1,q_2,Q)
\sum_{j'',m''}
\bra{j',m'}T^{(k_1)}_{q_1}\ket{j'',m''}
\bra{j'',m''}S^{(k_2)}_{q_2}\ket{j,m}.
\end{align*}
代入 Wigner--Eckart 定理：
\begin{align*}
&=\sum_{q_1,q_2,j'',m''}
C(k_1,k_2,K;q_1,q_2,Q)(-1)^{j'-m'+j''-m''}\\
&\quad\times
\begin{pmatrix}j'&k_1&j''\\-m'&q_1&m''\end{pmatrix}
\begin{pmatrix}j''&k_2&j\\-m''&q_2&m\end{pmatrix}\\
&\quad\times
\braket{j'\|T^{(k_1)}\|j''}
\braket{j''\|S^{(k_2)}\|j}.
\end{align*}
对 $m'',q_1,q_2$ 求和（利用 $q_1+q_2=Q$，$m''=m+q_2$，$m'=m+Q$ 等约束），得到的 $3j$ 符号求和恰好是 $6j$ 符号的定义（$3j$ 符号的 Racah 求和公式）：
\begin{align*}
&\sum_{m'',q_1,q_2}(-1)^{\cdots}
\begin{pmatrix}j'&k_1&j''\\-m'&q_1&m''\end{pmatrix}
\begin{pmatrix}j''&k_2&j\\-m''&q_2&m\end{pmatrix}\\
&\qquad\times
\begin{pmatrix}j'&K&j\\-m'&Q&m\end{pmatrix}\\
&\quad=(-1)^{\cdots}\sqrt{2K+1}\sqrt{2j''+1}
\begin{Bmatrix}k_1&k_2&K\\j&j'&j''\end{Bmatrix}
\begin{pmatrix}j'&K&j\\-m'&Q&m\end{pmatrix}.
\end{align*}
两边的 $3j$ 符号 $\begin{pmatrix}j'&K&j\\-m'&Q&m\end{pmatrix}$ 约去，即得 \cref{eq:reduced-me-product}。
\end{proof}
这个公式在原子物理中极为重要。例如，计算两个电子间的 Coulomb 相互作用 $1/r_{12}$（可展开为球谐函数的乘积）的矩阵元时，需要处理两个张量算符的耦合，$6j$ 符号自然出现并编码了不同耦合方案之间的变换。
\medskip
\noindent\emph{超精细结构的电四极相互作用。}
原子核的四极矩与电子云的相互作用为
\[
H_Q=-\frac{e^2Q}{2I(2I-1)}\sum_{q=-2}^{2}(-1)^q Q^{(2)}_q(\text{核})\,Q^{(2)}_{-q}(\text{电子}),
\]
这是两个秩 $2$ 张量算符的标量积。由 \cref{eq:scalar-product-me}，在 $\ket{F,m_F}$ 态中（$F$ 是总角动量 $\vb F=\vb I+\vb J$）：
\[
\bra{F,m_F}H_Q\ket{F,m_F}
=\frac{1}{2F+1}\braket{F\|Q^{(2)}(\text{核})\|F}\braket{F\|Q^{(2)}(\text{电子})\|F}.
\]
核约化矩阵元 $\braket{F\|Q^{(2)}(\text{核})\|F}$ 正比于核四极矩 $Q$，电子约化矩阵元通过 \cref{eq:reduced-me-product}（$k_1=2$ 作用于核空间，$k_2=2$ 作用于电子空间）分解为 $6j$ 符号和各自空间的约化矩阵元。最终能量修正为
\[
E_Q=\frac{B}{4}\left[\frac{3K(K+1)-4I(I+1)J(J+1)}{I(2I-1)J(2J-1)}\right],
\]
其中 $K=F(F+1)-I(I+1)-J(J+1)$，$B$ 是四极耦合常数。$6j$ 符号 $\begin{Bmatrix}I&J&F\\J&I&2\end{Bmatrix}$ 编码了核--电子角动量耦合的几何，而物理参数 $Q$ 和电子分布都进入常数 $B$。
\par\medskip
Wigner--Eckart 定理的本质是 Schur 引理在矩阵元层面的具体化。回顾 Schur 引理：若 $D$ 是不可约表示，则与所有 $D(g)$ 对易的矩阵是标量矩阵。在矩阵元语言中，这恰好说"不可约表示之间的交织算符只有零或标量倍"。
Wigner--Eckart 定理处理的是更一般的情形：$T^{(k)}$ 是从表示 $D^{(j)}$ 到 $D^{(j')}$ 的交织算符，中间经过 $D^{(k)}$ 的耦合。由 Schur 引理，$\operatorname{Hom}_G(D^{(j)},D^{(j')})$ 的维数是 $\delta_{jj'}$（不可约表示之间的交织算符），但 $T^{(k)}$ 不是直接的交织算符——它先与 $D^{(k)}$ 耦合再作用。展开 \cref{eq:tensor-op-expansion} 中，每个不可约表示 $D^{(J)}$ 通道的系数 $\alpha(J)$ 正是 $\operatorname{Hom}_G(D^{(J)},D^{(j')})$ 中的交织算符，由 Schur 引理 $\alpha(J)=0$ 除非 $J=j'$，且 $\alpha(j')$ 是标量。这就是约化矩阵元。
这个视角揭示了 Wigner--Eckart 定理的一般群论推广：对任意紧群 $G$，不可约张量算符的矩阵元正比于相应的 Clebsch--Gordan 系数（广义 $3j$ 符号），比例因子是约化矩阵元。对有限群，这正是\chapref{chap:group-theory-quantum}中对称性适配线性组合构造和投影算符方法的群论基础——投影算符 $\hat P^{(\Gamma)}_{ij}=\frac{d_\Gamma}{|G|}\sum_g D^{(\Gamma)*}_{ij}(g)\,g$ 正是"提取不可约表示 $\Gamma$ 的 $j$ 行 $i$ 列分量"的算符，其矩阵元由表示系数（有限群的"CG 系数"）完全确定，动力学信息只进入约化矩阵元。
本章到这里完成了从“几何转动”到“量子角动量耦合”再到“Wigner--Eckart 定理”的闭环。\chapref{chap:point-space-groups}与\chapref{chap:group-theory-quantum}中的点群方法解决离散空间对称性，\chapref{chap:rotation-group}解决完整转动对称性、自旋带来的二重覆盖，以及张量算符矩阵元的群论约化；下一章转向另一类完全不同但同样基础的离散对称性——全同粒子的置换对称性。
\section{习题与解答}
\label{sec:rotation-exercises}
以下习题检验角动量耦合、再耦合符号和张量算符的实际计算。习题本身不另设框，解答只做轻量区分。
\begin{enumerate}[label=\arabic*.,leftmargin=2.5em]
\item 从 $K(\uv n)\vb r=\uv n\times\vb r$ 出发证明 \eqnrefs{eq:K-identities}，并由此直接验证 Rodrigues 矩阵满足 $R^{\mathsf T}R=I_3$。
\emph{解答.} 向量三重积给出
\[
K^2\vb r=\uv n(\uv n\cdot\vb r)-\vb r,
\]
故 $K^2=\uv n\uv n^{\mathsf T}-I$。再乘一个 $K$，因 $K\uv n=0$，得 $K^3=-K$。又 $K^{\mathsf T}=-K$。令 $R=I+sK+(1-c)K^2$，其中 $s=\sin\theta,c=\cos\theta$，则
\[
R^{\mathsf T}=I-sK+(1-c)K^2.
\]
相乘并用 $K^4=-K^2$，$K$ 的奇次项抵消，$K^2$ 系数为
\[
2(1-c)-s^2-(1-c)^2=0,
\]
所以 $R^{\mathsf T}R=I$。
\item 证明 $\tr R_{\uv n}(\theta)=1+2\cos\theta$。
\emph{解答.} 由 Rodrigues 公式，$\tr K=0$，且
\[
\tr K^2=\tr(\uv n\uv n^{\mathsf T}-I)=1-3=-2.
\]
所以
\[
\tr R=3+(1-\cos\theta)(-2)=1+2\cos\theta.
\]
\item 用 Pauli 恒等式证明 $(\uv n\cdot\boldsymbol{\sigma})^2=I_2$，其中 $\lvert\uv n\rvert=1$。
\emph{解答.}
\[
(\uv n\cdot\boldsymbol{\sigma})^2
=n_i n_j(\delta_{ij}I+\ii\epsilon_{ijk}\sigma_k)
=\lvert\uv n\rvert^2I=I,
\]
因为 $n_i n_j$ 对 $i,j$ 对称，而 $\epsilon_{ijk}$ 反对称，第二项为零。
\item 证明 $U(\uv n,2\pi)=-I_2$、$U(\uv n,4\pi)=I_2$。
\emph{解答.} 直接代入 \eqnrefs{eq:su2-axis-angle}：
\[
U(2\pi)=\cos\pi I-\ii\sin\pi\,\uv n\cdot\boldsymbol{\sigma}=-I,
\]
而 $U(4\pi)=\cos2\pi I=I$。
\item 设 $U=a_0I-\ii\vb a\cdot\boldsymbol{\sigma}$、$V=b_0I-\ii\vb b\cdot\boldsymbol{\sigma}$。求 $UV$ 的同型参数，并说明 $SU(2)$ 的群乘法等价于单位四元数乘法。
\emph{解答.} 利用 Pauli 乘法，
\[
UV=(a_0b_0-\vb a\cdot\vb b)I
-\ii(a_0\vb b+b_0\vb a+\vb a\times\vb b)\cdot\boldsymbol{\sigma}.
\]
这正是单位四元数的乘法规则；若两组参数都满足 $a_0^2+a^2=b_0^2+b^2=1$，乘积也满足同一单位球面约束。
\item 由角动量代数推导 $[J^2,J_z]=0$ 和 $[J_z,J_\pm]=\pm\hbar J_\pm$。
\emph{解答.} 例如
\[
[J_z,J_+]=[J_z,J_x]+\ii[J_z,J_y]
=\ii\hbar J_y+\hbar J_x=\hbar J_+.
\]
另一个符号同理。对 $J^2$，
\[
[J_x^2,J_z]=J_x[J_x,J_z]+[J_x,J_z]J_x
=-\ii\hbar(J_xJ_y+J_yJ_x),
\]
而 $[J_y^2,J_z]$ 给出相反项，$[J_z^2,J_z]=0$，故总和为零。
\item 对 $j=1$ 写出 $J_z,J_+,J_-$ 的矩阵。
\emph{解答.} 在基 $\ket{1,1},\ket{1,0},\ket{1,-1}$ 中，
\[
J_z=\hbar\begin{pmatrix}1&0&0\\0&0&0\\0&0&-1\end{pmatrix},
\quad
J_+=\hbar\sqrt2
\begin{pmatrix}0&1&0\\0&0&1\\0&0&0\end{pmatrix},
\quad
J_-=J_+^\dagger.
\]
\item 说明为什么半整数 $j$ 不能给出 $SO(3)$ 的单值表示。
\emph{解答.} $SO(3)$ 中 $2\pi$ 转动就是单位元，因此任何群表示必须把它送到单位矩阵。但半整数 $j$ 中
\[
D^{(j)}(2\pi)=(-1)^{2j}I=-I,
\]
所以它不能因子化为 $SO(3)$ 表示，只能成为 $SU(2)$ 或双群表示。
\item 对 $\ell=1,s=\tfrac12$，用 \eqnrefs{eq:spin-orbit-energy} 求 $j=3/2$ 与 $j=1/2$ 的自旋--轨道能量。
\emph{解答.} $\ell(\ell+1)=2$，$s(s+1)=3/4$。对 $j=3/2$，$j(j+1)=15/4$，故
\[
E_{3/2}=\frac{\lambda\hbar^2}{2}\left(\frac{15}{4}-2-\frac34\right)
=\frac{\lambda\hbar^2}{2}.
\]
对 $j=1/2$，$j(j+1)=3/4$，故
\[
E_{1/2}=\frac{\lambda\hbar^2}{2}\left(\frac34-2-\frac34\right)
=-\lambda\hbar^2.
\]
\item 验证 Clebsch--Gordan 维数恒等式
\[
(2j_1+1)(2j_2+1)
=\sum_{j=|j_1-j_2|}^{j_1+j_2}(2j+1).
\]
\emph{解答.} 若 $j_1<j_2$，交换两个角动量的标签不会改变恒等式，因此只需写出 $j_1\ge j_2$ 的计算。右侧是首项 $2(j_1-j_2)+1$、末项 $2(j_1+j_2)+1$、项数 $2j_2+1$ 的等差数列，直接求和即得左侧。
\item 从最高权态出发求 $1\otimes\tfrac12$ 的全部耦合态。
\emph{解答.} 分解为 $\frac32\oplus\frac12$。最高态
\[
\ket{\tfrac32,\tfrac32}=\ket{1,1}\ket{\uparrow}.
\]
用 $J_-$ 得
\[
\ket{\tfrac32,\tfrac12}
=\sqrt{\frac23}\ket{1,0}\ket{\uparrow}
+\sqrt{\frac13}\ket{1,1}\ket{\downarrow}.
\]
与之正交的 $M=1/2$ 态为
\[
\ket{\tfrac12,\tfrac12}
=-\sqrt{\frac13}\ket{1,0}\ket{\uparrow}
+\sqrt{\frac23}\ket{1,1}\ket{\downarrow}.
\]
其余 $M<0$ 态由继续作用 $J_-$ 或时间反演得到。
\item 证明两个自旋 $1/2$ 的三重态在粒子交换下对称，而单态反对称。
\emph{解答.} 交换算符 $P_{12}$ 作用为 $P_{12}\ket{a}_1\ket{b}_2=\ket{b}_1\ket{a}_2$。因此 $\ket{\uparrow\uparrow}$、$\ket{\downarrow\downarrow}$ 不变，
\[
P_{12}\frac{\ket{\uparrow\downarrow}+\ket{\downarrow\uparrow}}{\sqrt2}
=+\frac{\ket{\uparrow\downarrow}+\ket{\downarrow\uparrow}}{\sqrt2},
\]
而带负号的组合取负。这一结果将在\chapref{chap:permutation-group}与费米统计结合。
\end{enumerate}
